\documentclass[11pt]{article}
\usepackage[T1]{fontenc}
\usepackage{microtype}
\usepackage[margin=1.05in]{geometry}
\usepackage{amsmath,amssymb,amsthm,booktabs,array}
\usepackage[hidelinks]{hyperref}

\newtheorem{theorem}{Theorem}
\newtheorem{lemma}[theorem]{Lemma}
\newtheorem{proposition}[theorem]{Proposition}
\newtheorem{corollary}[theorem]{Corollary}
\theoremstyle{remark}
\newtheorem{remark}[theorem]{Remark}
\newcommand{\Z}{\mathbb Z}
\newcommand{\Q}{\mathbb Q}
\newcommand{\R}{\mathbb R}
\newcommand{\C}{\mathbb C}
\DeclareMathOperator{\vp}{v_p}
\DeclareMathOperator{\Res}{Res}
\DeclareMathOperator{\lc}{lc}

\title{The Strong Factorial Conjecture for Univariate Binomials}
\author{Roy van Rijn}
\date{}
\hypersetup{pdftitle={The Strong Factorial Conjecture for Univariate Binomials},
  pdfauthor={Roy van Rijn}}

\begin{document}
\maketitle

\begin{abstract}
We prove that two consecutive positive factorial moments of a nonzero
univariate binomial cannot both vanish over a field of characteristic
zero. Equivalently, the two corresponding rational moment polynomials
are coprime for every support and every consecutive pair of orders.
A Pad\'e resultant bounds the orders at which a common factor can occur.
At the remaining orders, prime divisors of short factorial products
and the widths of shared Newton polygon edges exclude possible factor
degrees. An exact finite certificate completes this range.
\end{abstract}

\section{The problem and the proof}

Let \(K\) be a field of characteristic zero. The factorial functional on
\(K[X]\) is the \(K\)-linear map \(\mathcal L\) defined by
\(\mathcal L(X^j)=j!\). Edo and van den Essen's Strong Factorial
Conjecture asserts that a polynomial with \(r\) nonzero terms has a
nonzero factorial moment in every run of \(r\) consecutive positive
orders \cite{edo-essen}. We prove the case of two terms in one variable.

\begin{theorem}\label{thm:main}
For all integers \(a\ge0\), \(d\ge1\), \(N\ge2\) and every nonzero
\(f=X^a(A+BX^d)\in K[X]\), at least one of
\(\mathcal L(f^{N-1})\) and \(\mathcal L(f^N)\) is nonzero.
\end{theorem}

For a single monomial this follows from the definition. If \(AB\ne0\),
put \(t=B/A\) and define
\begin{equation}\label{eq:moment}
 P_m(t)=\sum_{k=0}^m\binom{m}{k}(am+1)_{dk}\,t^k,
 \qquad R_m(x)=x^m P_m(1/x),\qquad
 \mathcal L(f^m)=A^m(am)!\,P_m(B/A).
\end{equation}
Here \((y)_j=y(y+1)\cdots(y+j-1)\) and \((y)_0=1\).
The integral polynomials \(R_m\) are monic, with nonzero constant term.
It suffices to prove \(\gcd(R_{N-1},R_N)=1\) over \(\Q\),
since a rational B\'ezout identity holds over every \(K\).
Conversely, a nonconstant common factor gives a common root over
\(\C\) and hence two consecutive zero moments.

Rocks studied consecutive coprimality for two monomials and proved
several cases, using the same coefficient polynomials
\cite[Section 3.4]{rocks}. For the constant-plus-power support
\(\{0,d\}\), Filaseta and Rocks proved that, for each fixed \(d\), the
orders at which \(P_m\) is irreducible have lower density at least
\(\log 2\) \cite[Theorem 1.1]{filaseta-rocks}. Our conclusion concerns
every pair of consecutive orders and allows both exponents to vary.

Our proof uses prime divisibility at small orders and combines it with
a complex estimate at large orders. A Pad\'e truncation of
\(e^zR_{N-1}(-z^d)\), after normalization, has integral coefficients
and a long gap in its Taylor series. A Wronskian gives a nonzero
resultant with a factor lying above the roots of a putative common
factor. Primes dividing the factor's constant term give a lower
bound for the resultant by divisibility. The complex estimate,
averaged over rotated roots, gives an incompatible upper bound.
The remaining orders are finite. There the prime sieve and Newton
polygons leave only a short list of full polynomial gcd computations.

For an example of the prime sieve, take \(a=8\), \(d=3\), \(N=25\)
and \(p=101\). In both factorial intervals \((8m,11m]\),
\(m=24,25\), the only multiple of \(p\) is \(202\). All binomial
coefficients are \(p\)-adic units. The coefficient valuations of
\(P_{24}\) jump from zero to one at \(k=4\), and those of \(P_{25}\)
at \(k=1\). The lower Newton polygons give positive reciprocal
root valuations \(1/21\) and \(1/25\), respectively. Every root of
\(R_{25}\) has positive valuation, since \(202\) appears in the
first product of three factorial arguments. Thus the two orders
have no common root. In general, primes in the first few factorial
arguments constrain the constant term of a common factor. Newton
polygons then restrict the factor degrees allowed at those primes.

\section{Reduction of the exponent range}
\label{sec:supports}

Write \(P_m=P_{m,a,d}\) for the normalized moment polynomial and
\(R_m(x)=x^mP_m(1/x)\) for its monic reciprocal. Put \(b=a+d\) and
\(s=d^2/b\). For \(x\in\mathbb R\) and even \(m\), the Gamma integral gives
\[
 R_m(x)=\frac{1}{(am)!}
       \int_0^\infty u^{am}(x+u^d)^m e^{-u}\,du>0 .
\]
This observation will also be used in the arithmetic argument.

\begin{lemma}[Root annuli and separation of the shapes]
\label{lem:gamma-annulus}
For \(a\ge1\), every root \(x\) of \(R_m\) satisfies
\[
 am<|x|^{1/d}\le (a+d)m .
\]
For any \(a\ge0\), the pair \(P_{N-1},P_N\) is coprime if
\(a\ge d(N-2)\). If \(N\ge25\), it is also coprime if
\(6a\ge (5N-11)d\).
\end{lemma}

The proof is in Appendix~\ref{sec:gamma-annulus-proof}. It uses the
root annuli of the unit-step Gamma polynomials and the fact that
these annuli multiply under composition in the binomial basis.
The strengthened inner radius matters because the factorial arguments
in the two moments grow at different rates.

\begin{lemma}[Constant plus a power]
\label{lem:constant-support}
For \(a=0\) and every \(d\ge1\), the consecutive polynomials
\(P_{N-1,0,d},P_{N,0,d}\) are coprime over \(\mathbb Q\).
\end{lemma}

\begin{proof}
In this case \(R_N'=NR_{N-1}\), by coefficient comparison. The
squarefreeness proof in Appendix~\ref{app:constant-proof} therefore
gives the result.
\end{proof}

\begin{lemma}[Endpoint normalization]
\label{lem:center-ratio}
For \(a\ge1\) define
\[
 G_m=\left(\frac{((a+d)m)!}{(am)!}\right)^{1/m},
 \qquad
 P_m(w/G_m)=1+w^m+\sum_{k=1}^{m-1}c_{m,k}w^k .
\]
For \(1\le k<m\) one has
\begin{equation}\label{eq:normalized-coefficients}
 0<c_{m,k}\le
 \binom{m}{k}\exp\!\left(-\frac{s\,k(m-k)}{2m}\right).
\end{equation}
For \(N\ge2\) the centers satisfy
\[
 G_N/G_{N-1}>(N/(N-1))^{d/2}.
\]
\end{lemma}

\begin{proof}
The consecutive differences of \(\log((am+dk)!)\) are the logarithms
of rising products of length \(d\). For \(1\le k<m\) the second
difference is
\[
 \sum_{j=1}^d\log\!\left(1+\frac{d}{am+d(k-1)+j}\right)
 \ \ge\ \frac{d^2}{(a+d)m}.
\]
Indeed, \(\log(1+v)\ge v/(1+v)\) for \(v\ge0\). Subtracting the
quadratic with this second difference and comparing with the chord
through the endpoints proves \eqref{eq:normalized-coefficients}.

For the second claim, let \(f(u)=\log\Gamma(u+1)\) and
\(H_0(x)=(f((a+d)x)-f(ax))/x\). The trigamma series gives
\(f''(u)>1/(u+1)\), by comparison with the integral. Since
\(a x\ge1\) for \(x\ge1\),
\[
 H_0'(x)=\frac{1}{x^2}\int_{ax}^{(a+d)x}u f''(u)\,du
       >\frac{d}{2x}.
\]
Integrate from \(N-1\) to \(N\).
\end{proof}

\begin{lemma}[Ratio bound]
\label{lem:ratio-range}
For \(a\ge1,d\ge1,N\ge25\), if \(s\ge\sigma_N\), where
\[
 \sigma_N=
 \begin{cases}
  2\log N+3& (25\le N<200),\\
  2\log N+2& (N\ge200),
 \end{cases}
\]
then \(P_{N-1},P_N\) are coprime.
\end{lemma}

\begin{proof}
Suppose \(s\ge\sigma_N\). The coefficient bound in
\eqref{eq:normalized-coefficients} is symmetric in \(k,m-k\).
Write \(\sigma_N=2\log N+C\) and take \(m=N-1\) or \(N\).
The following rational constants bound the short and long parts
of the coefficient sum.
\[
\begin{array}{c|cccc}
 N&A&E&q&\eta\\ \hline
 25\le N<200&.224&1.218&.52&7/8\\
 N\ge200&.368&1.033&.24&19/20
\end{array}
\]
For \(1\le k\le4\), use \(\binom mk\le N^k/k!\) in
\eqref{eq:normalized-coefficients}. The factor \(N^k\) cancels
\(\exp(-k\log N)\). The resulting bound is \((AE^k)^k/k!\), since
\[
 A\ge e^{-C/2},\qquad
 E\ge \exp\!\left(\frac{\log N+C/2}{N-1}\right).
\]
For \(5\le k\le m/2\), put \(u=k/m\) and use
\(\binom mk\le(e/u)^k\). The logarithm of the resulting bound,
divided by \(k\), is \(1-\log u-\sigma_N(1-u)/2\).
Its second derivative in \(u\) is \(1/u^2>0\), so it is bounded
above by the larger value at \(u=5/m\) and \(u=1/2\).
The corresponding exponentials are bounded by \(eAE^5/5\)
and \(2e\,e^{-\sigma_N/4}\).
On a row with fixed \(C\), the quotient \((\log N+C/2)/(N-1)\)
decreases because its derivative has numerator
\(1-1/N-\log N-C/2<0\). The second endpoint exponential
decreases because \(\sigma_N\) increases.
Thus it suffices to check the bounds on \(E\) and \(q\) at
\(N=25\) and \(N=200\), respectively.

The two rows satisfy
\[
 2\sum_{k=1}^{4}\frac{(AE^k)^k}{k!}+\frac{2q^5}{1-q}<\eta .
\]
The supplement checks this rational inequality and encloses the
exponentials at the two endpoints. By symmetry the displayed
sum bounds \(\sum_{k=1}^{m-1}c_{m,k}\), including the middle
coefficient when \(m\) is even. At a root \(w\) of \(P_m(w/G_m)\) with \(|w|\le1\),
\(1\le |w|^m+\sum_{k=1}^{m-1}c_{m,k}<|w|^m+\eta\).
Apply the same argument to \(1/w\) in the reversed polynomial
when \(|w|\ge1\). Every root in the normalized variable then
satisfies
\[
 (1-\eta)^{1/m}<|w|<(1-\eta)^{-1/m}.
\]

At a common root the two values of \(|w|\) have ratio \(G_N/G_{N-1}\).
By Lemma~\ref{lem:center-ratio} this is too large for the two annuli, since
\[
 \frac d2\log\!\left(\frac N{N-1}\right)
              \frac{N(N-1)}{2N-1}
 \ \ge\ \frac{d(N-1)}{2(2N-1)}
 \ >\ -\log(1-\eta).
\]
For the last inequality, \(d>s\ge\sigma_N\) gives \(d\ge10\)
in the first row and \(d\ge13\) in the second. The minimum
of the rational lower bound in each row occurs at the smallest
\(N\) and the smallest \(d\). The resulting bounds are \(120/49\)
and \(2587/798\), which exceed \(\log 8\) and \(\log 20\),
respectively. The supplement encloses both logarithms and
checks the lower bounds on \(d\). This contradicts a common root.
\end{proof}


\begin{lemma}[Support reduction]
\label{lem:support-reduction}
If \(N\ge25\) and \(P_{N-1,a,d},P_{N,a,d}\) have a common root, then
\[
 a\ge5,\qquad 2\le d<3a,\qquad
 a<d(N-2),\qquad 6a<(5N-11)d,\qquad
 s<\sigma_N,
\]
where \(s=d^2/(a+d)\).
\end{lemma}

\begin{proof}
Lemma~\ref{lem:constant-support} disposes of \(a=0\).
Lemma~\ref{lem:positive-boundary} covers \(a=1,d\ge3\) and
\(a\ge2,d\ge3a\). The remaining small positive
lower exponents and the adjacent exponents are dealt with in
Proposition~\ref{prop:additional-boundary}. These exclusions give
\(a\ge5\) and \(2\le d<3a\). Lemma~\ref{lem:gamma-annulus} gives
the two bounds on the shape. Lemma~\ref{lem:ratio-range} gives the
bound on \(s\).
\end{proof}

\subsection{Terminals through twenty-four}

\begin{proposition}
\label{prop:small-terminals}
Theorem~\ref{thm:main} holds for \(2\le N\le24\).
\end{proposition}
\begin{proof}
Lemma~\ref{lem:constant-support} covers \(a=0\). At \(N=2\), \(R_1\)
has a real root and \(R_2\) is positive on \(\R\) by its gamma integral.
Take \(a\ge1\) and \(3\le N\le24\). The weak annulus in
Lemma~\ref{lem:gamma-annulus} gives \(a<d(N-2)\) if the two polynomials
share a root.

Write \(s=d^2/(a+d)\) and use the endpoint normalization
of Lemma~\ref{lem:center-ratio}. If \(s\ge20\), the mixed
coefficients at either order have sum less than \(1/2\).
Indeed their upper bound in \eqref{eq:normalized-coefficients}
is symmetric, and for \(k\le m/2\) the exponential term is
at most \(e^{-5k}\). Thus the sum is bounded by
\[
 2\bigl((1+e^{-5})^{24}-1\bigr)
 <2\bigl(e^{6/35}-1\bigr)<\frac{12}{29}<\frac12.
\]
Here \(e^{-5}<1/140\), and \(e^x<1/(1-x)\) for \(0<x<1\).
In the normalized variable the roots have moduli between
\(2^{-1/m}\) and \(2^{1/m}\), by comparison of \(1+w^m\)
with the mixed terms and by reversal. Since \(d>s\ge20\),
\[
 \frac d2\log\!\left(\frac{N}{N-1}\right)
 \frac{N(N-1)}{2N-1}
 \ge \frac{d(N-1)}{2(2N-1)}\ge 4>\log 2 .
\]
The center ratio in Lemma~\ref{lem:center-ratio} therefore
separates the annuli. A common root requires \(s<20\).

We next show that \(d<43\). If \(d\ge43\), set
\(r=\lfloor d/43\rfloor\). Then \(43\le d/r<86\), and the interval
\((aN/r,(aN+2d)/r]\) has length at least \(86\). Its upper endpoint
is below \(86[N(N-2)+2]\le45580\).
The bound \(s<20\) gives \(a>d^2/20-d\ge d^2/43\ge rd\), so the
left endpoint exceeds \(dN\). The first record prime gap larger
than \(86\) starts at \(360653\) \cite[Table~8]{oliveira-gaps}.
It follows that the interval contains a prime \(p>dN\), and that
\(rp\) belongs to the first \(2d\) factorial terms at order \(N\).

Since \(a<d(N-2)\), we have \(p^2>(a+d)N\).
Each factorial interval \((am,(a+d)m]\) contains at most one
multiple \(u_m\) of \(p\), and all binomial coefficients are
\(p\)-adic units. A positive valuation of a reciprocal root
has the form \(1/w_m\), where
\(w_m=1+\lfloor((a+d)m-u_m)/d\rfloor\).
Equal positive valuations at both orders would imply
\(u_N-u_{N-1}\in[a+1,a+2d-1]\subset(0,p)\), which is impossible.
Thus all roots of a common factor would be \(p\)-adic units.
But the first \(2d\) terms at order \(N\) contain \(rp\), so
\(R_N\bmod p\) has at most one root away from zero, counted with
multiplicity. A monic common divisor reduces with its degree
unchanged, and would have nonzero constant term and even degree
at least two. This contradiction proves \(d<43\).

It remains to test the integer triples
\[
 3\le N\le24,\qquad 1\le d\le42,\qquad
 1\le a<d(N-2),\qquad d^2<20(a+d).
\]
The supplementary verifier exhausts this set and finds
\(217\,952\) triples. At each triple it forms the two monic
reciprocals modulo \(1\,000\,000\,007\), a prime certified by
trial division, and computes their gcd by Euclidean division.
Every gcd is one. A monic integral common divisor would remain
a nonconstant common divisor on reduction. This proves the proposition.
\end{proof}

\section{Prime bounds and the finite range}
\label{sec:arithmetic}

Take \(a,d\ge1\), \(N\ge25\), \(d<3a\), \(6a<(5N-11)d\), and
\(s=d^2/(a+d)<\sigma_N\), as in Lemma~\ref{lem:support-reduction}.
In this section \(m\) runs through \(N-1,N\), \(b=a+d\), \(M=bN\),
and \(F\) is a nonconstant monic rational common divisor of
\(R_m(x)=x^mP_m(1/x)\). By Gauss's lemma \(F\in\Z[x]\), and \(F(0)\ne0\).
Write \(h=\deg F\) and \(H=dh\). The reciprocal of even order is
positive on \(\R\), as is seen by integrating \((x+u^d)^m\) against
the gamma density of shape \(am+1\). Hence \(h\) is even.

In fact \(h<N-1\). To see this, put \(n=N-1\) and \(u=an+1\).
If \(P_n\mid P_{n+1}\), the quotient is \(1+ct\). Its linear
coefficient would have to be both
\[
 c_1=(n+1)(u+a)_d-n(u)_d,\qquad
 c_2=\frac{(u+dn)_b}{(u)_a}.
\]
Set \(A=(u)_b\), \(B=(u)_a(u)_d\), and \(f(z)=(u+z)_b\).
Then \((u)_a(c_2-c_1)=f(dn)-(n+1)A+nB\).
Convexity gives \(f(dn)\ge nf(d)-(n-1)A\), and
\(f(d)/A=(u+b)_d/(u)_d>A/B\). Thus
\(f(d)-2A+B>(A-B)^2/B\ge0\), proving \(c_2>c_1\).

A sufficiently large prime can occur in a common root valuation only
through a multiple close to the end of each factorial product. The
locations of those multiples must agree to within one coefficient step
after translation by \(b\). We first use this observation to restrict
primes in the two initial factorial blocks.

For a prime \(p\) put \(y_m(k)=v_p([t^k]P_m)\) and
\[
 U_m(j)=y_m(m)-y_m(m-j)
       =v_p\!\left(\prod_{i=0}^{dj-1}(bm-i)\right)-v_p\binom mj .
\]
A lower Newton slope \(\lambda\) of \(P_m\) is a valuation of roots of
\(R_m\). Its width is the number of roots at that valuation, counted
with multiplicity. Equivalently the extremes of the maximizing indices
of \(U_m(j)-\lambda j\) have this difference.

\begin{lemma}[Degree allocations at a prime]
\label{lem:newton-alloc}
For a common slope with reduced denominator \(v\), let \(w\) be the
smaller of its two widths. The common factor can use only a number of
roots from \(\{0,v,2v,\ldots,w\}\) at that slope. Consequently \(h\)
belongs to the sum of these sets over the common slopes, at every prime.
In particular each positive valuation of a root of \(F\) is at least
\(1/h\).
\end{lemma}
\begin{proof}
At each slope the roots of a common factor occur with at most the
smaller multiplicity. Its Newton polygon has integral vertices, so
each width in that polygon is divisible by the reduced denominator
of the slope. The last assertion follows since such a width is at
most \(h\).
\end{proof}

\begin{lemma}[Prime support of short blocks]
\label{lem:prime-support}
Set
\[
 C_N=\max\{a+2d-1,\sqrt M\},\qquad B_h=\min\{dN,\max(2H,C_N)\}.
\]
Every prime divisor of \(F(0)\) or of either block
\(\prod_{i=1}^H(am+i)\) is at most \(B_h\).
\end{lemma}
\begin{proof}
If \(p>dN\), then \(p^2>M\) because \(a<d(N-2)\), and all binomial
coefficients are \(p\)-adic units. Each full factorial interval
\((am,bm]\) has at most one multiple \(u_m\) of \(p\). If present,
it yields only one positive slope, equal to \(1/w_m\) with
\(w_m=1+\lfloor(bm-u_m)/d\rfloor\).
A common positive slope would require the same \(w_m\) in both orders.
It would give
\begin{equation}
 u_N-u_{N-1}\in[b-d+1,b+d-1]=[a+1,a+2d-1].
 \label{eq:translated-multiples}
\end{equation}
This interval is positive and below \(p\), so it contains no multiple
of \(p\).

Now let \(p>\max(2H,C_N)\) and suppose that \(F\) has a positive slope
\(\lambda\) at \(p\). By Lemma~\ref{lem:newton-alloc},
\(\lambda\ge1/h\). Since \(p^2>M\), the count of terminal multiples
gives \(U_m(j)<1+dj/p\) for \(j>0\). The maximizing value of
\(U_m(j)-\lambda j\) is at least zero, so every maximizing index is
less than \(2h\). Throughout this range \(j<p\) and \(dj<p\).
The binomial valuation cannot decrease because \(j!\) is a unit.
The factorial valuation can rise only once and by one. The weighted
ordinate therefore decreases strictly between jumps. A positive
slope must tie zero to the index of the single factorial jump, giving
slope \(1/w_m\) with \(w_m=1+\lfloor(bm-u_m)/d\rfloor\).
As before, a shared slope implies \eqref{eq:translated-multiples},
which is impossible since \(p>C_N\).

These two cases show that \(p>B_h\) cannot divide \(F(0)\).
If it divides one of the first \(H\) factorial terms, all the
coefficients at \(k\ge h\) in \(P_m\) vanish modulo \(p\).
Then \(R_m\bmod p\) is a power of \(x\) times a polynomial of degree
at most \(h-1\). Its monic divisor \(F\bmod p\) has degree \(h\)
and nonzero constant term, a contradiction.
\end{proof}

\subsection{The degree bound}

\begin{lemma}[Degree bound from the first factorial block]
\label{lem:degree-cap}
If \(aN+1>\max\{H+13,279H/262\}\), then \(2H<C_N\).
The starting-point inequality holds whenever \(a\ge2d\), or
whenever \(N\ge275\) and \(aN\ge1375\). In particular the second
condition covers \(a\ge5\) at these terminals and every \(a\ge1\)
when \(N\ge3150\).
\end{lemma}

\begin{proof}
Theorem~1(a) of Laishram and Shorey~\cite{laishram-shorey} gives a prime
divisor greater than \(2H\) of the first \(H\) consecutive integers
beginning at \(aN+1\), under the stated starting-point inequality.
Lemma~\ref{lem:prime-support} then gives \(2H<C_N\).
If \(a\ge2d\), the inequality follows from \(h\le N-1\) and \(N\ge25\).

It remains to establish the starting-point inequality for
\(N\ge275\), \(aN\ge1375\), and \(a<2d\).
Put \(X=aN\), \(w=X/40\), and consider
three intervals of length \(w\) whose left endpoints are
\[
 X,\qquad X+2w+4d,\qquad X+6w+12d.
\]
Take each interval open at its left endpoint. Their union is contained
in \((X,X+E]\) with \(E=7w+12d< X/3\), since \(d<3a\).
Each interval contains a prime. Indeed \(X\ge1375\). For
\(X\le7160\), \(7160<X\le11750\), and \(11750<X<89693\), respectively,
the upper endpoint \(4X/3\) lies below \(9551,15683,155921\),
where the record prime gaps first exceed \(34,36,72\),
respectively~\cite[Table~8]{oliveira-gaps}. In each range \(w\)
exceeds the corresponding gap. For \(X\ge89693\),
Dusart~\cite[Proposition~5.4]{dusart-interval} gives a prime within
\(x/\log^3x<w\) of each left endpoint \(x\in[X,4X/3)\).

Discard any interval meeting
\[
 I_j=[(M-b)/j,(M+2d)/j]\qquad(1\le j\le4).
\]
The envelope misses \(I_1\), since \(M-b>(3/2)(1-1/N)X>4X/3\).
If \(I_j\), \(2\le j\le4\), meets the envelope, then \(M+2d>jX\).
For \(j'<j\) this implies
\[
 (M-b)/j' > X+(X-a-3d)/3>X+E,
\]
where \(E/X<7/40+36/275 < (1-10/275)/3\).
Thus at most one \(I_j\) meets the envelope. Its width is
at most \((a+3d)/2<5d/2\), smaller than the separation of the
three intervals. We obtain two primes \(p_1<p_2\) in distinct
retained intervals.

For either retained prime \(p\), we have \(p>X\ge m\) and
\(p^2>M\) because \(b<4a\). Therefore
\(y_m(k)=\lfloor(am+dk)/p\rfloor\). The avoidance of the \(I_j\)
gives the same \(q=\lfloor bm/p\rfloor\in\{1,2,3\}\) at both orders
and \((q+1)p-bm>2d\). The right endpoints of the constant-ordinate
intervals before the jumps are
\[
 k_j(m)=\left\lfloor\frac{jp-am-1}{d}\right\rfloor
 \qquad(1\le j\le q).
\]
Every positive edge before the last of these endpoints joins two such
endpoints and has slope at most \(1/\lfloor p/d\rfloor\).
This follows by subtracting their floor expressions, since an edge
of height \(J\) then has width at least \(J\lfloor p/d\rfloor\).
The edge from \(k_q(m)\) to \(m\) has height one and width
\[
 w_m=\left\lceil(bm-qp+1)/d\right\rceil
 \le \lceil p/d\rceil-2 <\lfloor p/d\rfloor.
\]
All remaining points lie on or above this final segment. Its slope
is larger than all earlier slopes, and \(w_N>w_{N-1}\) since \(b>d\).
Neither final slope is shared.

For an interior positive edge of height \(J_e\) in \(P_N\), its
width is \(w_e=J_e p/d+\epsilon_e\), with
\(|\epsilon_e|<1\). If the factor uses \(r_e\) of its roots,
the quantity \(r_eJ_e/w_e\) is an integer by
Lemma~\ref{lem:newton-alloc}. Summing over these edges yields
\[
 h=h_0+Jp/d+\epsilon,\qquad
 0\le h_0<(p-X)/d,\qquad J\in\{0,1,2\}.
\]
Here \(h_0\) counts slope-zero roots and \(J\) is at most the total
interior height \(q-1\). The error satisfies \(|\epsilon|<2\) when
\(J>0\), and is zero when \(J=0\).
For \(p_1,p_2\), the intervals of possible \(h\) with equal positive
\(J\) are disjoint because \(p_2-p_1>p_1-X+4d\), by the placement of
the intervals. Intervals for unequal \(J\) are disjoint as well,
since \(X>2E+4d\). Thus \(J=0\) at both primes. The smaller prime
comes from the first or second interval, giving
\[
 H<p_1-X\le 3w+4d<.12X .
\]
As \(X\ge1375\), this implies the required starting-point inequality.
\end{proof}

\subsection{Bounds from records of prime gaps}

For an even record gap \(G\), write \(X_G\) for the lower prime
endpoint of the first larger record gap in
\cite[Table~8]{oliveira-gaps}. Every real interval \((x,y]\)
with \(x\ge2\), \(y<X_G\) and \(y-x\ge G\) contains a prime.
Otherwise the consecutive primes straddling \(x\) have gap
greater than \(G\) and lower endpoint below \(X_G\).

\begin{lemma}[Uniform gap and degree bounds]
\label{lem:gap-caps}
Let \(G\ge52\) be an even record gap such that
\begin{equation}
 G\ge 4\sigma_N,\qquad
 G\left(\frac{N(5N-11)}6+2\right)<X_G .
 \label{eq:gap-conditions}
\end{equation}
Then \(d<G/2\). If Lemma~\ref{lem:degree-cap} applies, then also
\(H<G\).
\end{lemma}
\begin{proof}
Put \(g=G/2\). If \(d\ge g\), take \(r=\lfloor d/g\rfloor\).
Then \(g\le d/r<2g\), so \((aN/r,(aN+2d)/r]\) has length
at least \(G\). Its upper endpoint is below \(X_G\) by
\(6a<(5N-11)d\) and \eqref{eq:gap-conditions}. Its left endpoint
exceeds \(dN\), since \(s<\sigma_N\) and \(d\ge g\ge2\sigma_N\)
give \(a>d^2/\sigma_N-d\ge d^2/g\ge rd\).
A prime \(p\) in the interval divides \(rp\) among the first
\(2d\) terms of the upper factorial block, contrary to
Lemma~\ref{lem:prime-support}.

Suppose now that \(2H<C_N\) and \(H\ge G\).
Take \(r=\lfloor H/G\rfloor\). The interval
\((aN/r,(aN+H)/r]\) has length at least \(G\), and its upper
endpoint is less than
\[
 2G\left(\frac{N(5N-11)}{6h}+1\right)
 \le G\left(\frac{N(5N-11)}6+2\right)<X_G .
\]
The left endpoint exceeds two since \(r\le H/G<dN/G\) and
\(aG/d>G/3\).
If \(C_N=\sqrt M\), then \(aN/r\ge aGN/H>2aG\sqrt M/b>\sqrt M\).
If \(C_N=a+2d-1\), then \(2H<C_N\) and \(h\ge2\) give \(a>2d\).
Since \(d<G/2\), we have \(r<N/2\) and \(aN/r>2a>C_N\).
Thus in both cases the interval supplies a prime \(p>C_N\)
dividing one of the first \(H\) terms. This contradicts
Lemma~\ref{lem:prime-support}.
\end{proof}

Table~\ref{tab:gap-bounds} gives choices of \(G\) for
Lemma~\ref{lem:gap-caps}. Use the first applicable row. The
inequalities in \eqref{eq:gap-conditions} follow by checking
the right endpoint of each row, since \(2\log N+3\) is an
increasing upper bound for \(\sigma_N\).
\begin{table}[ht]
\centering
\begin{tabular}{r r r}
\hline
Upper terminal \(N\) & \(G\) & \(X_G\)\\
\hline
50 & 72 & \(155\,921\)\\
100 & 114 & \(1\,349\,533\)\\
300 & 154 & \(17\,051\,707\)\\
\(1\,000\) & 222 & \(189\,695\,659\)\\
\(2\,000\) & 282 & \(1\,294\,268\,491\)\\
\(6\,000\) & 354 & \(10\,726\,904\,659\)\\
\(10\,000\) & 456 & \(42\,652\,618\,343\)\\
\(19\,999\) & 468 & \(182\,226\,896\,239\)\\
\(32\,767\) & 516 & \(461\,690\,510\,011\)\\
\(65\,535\) & 602 & \(2\,614\,941\,710\,599\)\\
\(3\,269\,017\) & 1132 & \(43\,841\,547\,845\,541\,059\)\\
\hline
\end{tabular}
\caption{Prime-gap bounds used to control \(d\) and \(dh\). The record
endpoints are from \cite[Table~8]{oliveira-gaps}.}
\label{tab:gap-bounds}
\end{table}

\begin{lemma}[Refinement using the actual factorial position]
\label{lem:local-gap}
Assume \(2H<C_N\) and \(H\le H_0\). Let \(G\ge52\) be an even record
gap. If \(H_0/G<N/2\) and \(aN+H_0<X_G\), then
\(H<G\) and \(d<G/2\).
\end{lemma}
\begin{proof}
If \(H\ge G\), the interval with \(r=\lfloor H/G\rfloor\) in the
second part of Lemma~\ref{lem:gap-caps} has upper endpoint below
\(X_G\) by assumption, as well as length at least \(G\). The two
comparisons with \(C_N\) in that proof still hold, using
\(r\le H_0/G<N/2\) in the second case. The resulting prime gives
the same contradiction.
\end{proof}

In particular \(H\le281\) and \(d\le140\) for
\(1001\le N\le3295\) whenever Lemma~\ref{lem:degree-cap} applies.
Starting with \(G=354\) from
Table~\ref{tab:gap-bounds}, apply Lemma~\ref{lem:local-gap}
successively with
\[
\begin{array}{r|r|r}
G&\text{upper terminal }N&X_G\\ \hline
320&5000&3\,842\,610\,773\\
292&4150&2\,300\,942\,549\\
288&3467&1\,453\,168\,141\\
282&3295&1\,294\,268\,491
\end{array}
\]
At each step \(aN+H_0<(5/6)D_0N^2+H_0<X_G\), where \(d\le D_0\)
and \(H\le H_0\) are the preceding bounds.

\subsection{Exact computation in the finite range}

\begin{proposition}
\label{prop:finite}
Under the support reduction of Lemma~\ref{lem:support-reduction},
the consecutive polynomials are coprime for \(a\ge5\), \(d\ge2\),
and \(25\le N<3150\).
\end{proposition}
\begin{proof}
Use \(d^2<19b\), which follows from \(\sigma_N<19\) in this
range, and use the gap caps
\[
\begin{array}{c|ccccc}
N&25\text{--}50&51\text{--}100&101\text{--}300&301\text{--}1000&
1001\text{--}3149\\
d\le&35&56&76&110&140 .
\end{array}
\]
These come from Lemmas~\ref{lem:gap-caps} and \ref{lem:local-gap}.
For \(2200\le N<3150\), the analytic exclusions in
Appendix~\ref{app:middle-tail} also require \(b>N\) and
\(97d<3a\).

For each \(N,a\), the inequalities \(6a<(5N-11)d\), \(2\le d<3a\),
the displayed cap, and the two additional inequalities when
\(N\ge2200\) give an integer interval \([d_0,d_1]\) for \(d\).
When \(N\ge275\) or \(a\ge2d\), Lemma~\ref{lem:degree-cap}
reduces the bound in Lemma~\ref{lem:prime-support} to
\(\min(dN,C_N)\), independent of \(h\). Otherwise use \(dN\).
The resulting bound is nondecreasing in \(d\), including where
it switches to \(dN\). So a prime exceeding the bound at \(d_1\)
in either first \(2d_0\)-term block eliminates the whole interval.
The remaining gaps are tested individually with their first
\(2d\) terms and the conditions \(a<d(N-2)\), \(d^2<19b\).

For each remaining support, try the possible even degrees
\(2\le h<N-1\). Test the starting-point inequality directly
when \(N<275\). Impose the degree bound when it applies, and
test the full first \(H\) terms with
Lemma~\ref{lem:prime-support} in both orders.

The remaining degrees are tested at odd primes \(aN<p\le M\).
Here \(p>m,d\), \(p^2>M\), and \(am<p\). The ordinate is
\(\lfloor(am+dk)/p\rfloor\) and \(q_m=\lfloor bm/p\rfloor\le3\).
Thus the lower hull uses at most five points
\[
 (0,0),\qquad
 \left(\left\lfloor\frac{jp-am-1}{d}\right\rfloor,j-1\right)
 \quad(1\le j\le q_m),\qquad (m,q_m).
\]
Every omitted point is to the left of a retained point of the
same ordinate and lies on or above the retained lower hull,
which is nondecreasing. Compute the allowed degree sum from
Lemma~\ref{lem:newton-alloc} at each prime and intersect it with
the surviving degrees. For example \(N=25,a=5,d=8,p=127\)
gives hull slopes \(1/16,1/8\) and \(1/16,1/9\) at orders
24 and 25, respectively. Only degrees zero and sixteen are
allowed at this prime.

The supplementary verifier exhausts these integer intervals.
The largest factorial argument is \(1\,156\,520\,634\).
A segmented sieve stores smallest odd prime factors through that limit.
Trial divisions from this table and cross multiplication of slopes
perform the tests above. In prime caps the square-root threshold
is replaced by its integer part, checked by squaring. The initial prime-cap test leaves
\(21\,583\) groups \((N,a)\). The individual gap test leaves
\(74\,362\) supports. The degree and prime-block tests exclude
\(62\,593\), and the Newton allocations exclude \(11\,705\).
For the remaining \(64\) supports, the verifier forms both
monic reciprocals modulo the prime \(1\,000\,000\,007\) and
checks by Euclidean division that their gcd is one. The
modulus is checked by trial division. A monic integral common
divisor would retain its degree on reduction, so these
computations prove rational coprimality.
\end{proof}

\section{An exponential Padé estimate}\label{sec:pade}

Let \(F\) be a monic irreducible common factor of the reciprocal
moment polynomials \(R_{N-1},R_N\). Write \(h=\deg F\), \(H=dh\),
\(b=a+d\), \(M=bN\), \(m=N-1\), \(s=d^2/b\), and \(\alpha=d/b\).
We use the root bounds \(ai<|x|^{1/d}\le bi\) for \(R_i(x)=0\)
with \(i=m,N\), from Lemma~\ref{lem:gamma-annulus}. In particular,
\(H\log(aN)<\log|F(0)|<H\log M\). The Padé construction works on
\(z^d=-x\), where the root bound is linear in the moment order.
Its key feature is that an integral polynomial approximates the
same polynomial to high order after multiplication by each of the
\(d\) exponentials \(e^{-\omega z}\), \(\omega^d=1\).
The resulting Wronskian makes at least one of the approximants
nonzero above every root of \(F\). The comparison of an integer
resultant with its local estimates adapts the Laguerre construction
in Nguyen-Dang's preprint \cite[Proposition~3.9]{nguyen-dang}.
The following proof derives the rotated truncations, their
Wronskian and both local estimates for the binomial moments.

\begin{lemma}[Padé resultant inequality]\label{lem:resultant}
Assume \(N\ge200\), \(a,d\ge1\) and \(d<3a\). Set
\(r=(b-1)m\), and put
\[
 \begin{gathered}
 K_\pm=\tfrac12\pm\tfrac{\alpha}{14},\qquad
 \rho_0=b(1-\alpha K_+)=a+dK_-,\\
 g=\max\{b\log(b/\rho_0)-b+\rho_0,\,
          \rho_0-a-a\log(\rho_0/a)\}.
 \end{gathered}
\]
There is a monic irreducible factor \(G\) of \(F(-z^d)\) of degree
\(D=hk\), where \(1\le k\le d\), and \(K=|G(0)|\) satisfies
\(K^d=|F(0)|^k\). Write \(V=\log(M+1)\).
At each complex root \(z\) of \(G\) one may take the gain
\[
 L(z)=\max\left\{(m!e^{-mg-1})^{1/(m+1)},
                    \frac{am-\Re z}{d}\right\}.
\]
With \(\Lambda=(\prod_{G(z)=0}L(z))^{1/D}\), one has
\begin{equation}\label{eq:pade-norm}
 \log\Lambda<
 \sum_{p\mid K}\min\left\{\frac{\vp(K)}{D},\frac1{p-1}\right\}\log p
 +\frac{\omega(K)V}{N}.
\end{equation}
Here \(\omega(K)\) is the number of distinct prime divisors.
In addition, the first entry in \(L(z)\) can use \(g=0\)
when \(\Re z\le am\). For \(N\ge200\) the first entry is
greater than \(me^{-1.01-g}\). For \(d=1\) the additional gain
\(m\sqrt a\,e^{-1/2}/(1+a^{-1/2})\) is available.
\end{lemma}

\begin{proof}
Let \(B(z)=(am)!R_m(-z^d)/m!\) and let \(Q(z)\) be the Taylor
polynomial of \(e^z B(z)\) through degree \(r\). For \(0\le j\le bm\)
and \(v=bm-j\), direct expansion gives
\begin{equation}\label{eq:pade-coeff}
 [z^j]e^zB(z)=\frac{v!}{m!}
 [w^v](1+w)^{am}((1+w)^d-1)^m.
\end{equation}
Indeed the term indexed by \(i\) in \(R_m\) contributes
\((-1)^{m-i}\binom mi(am+di)!/(m!(am+di-v)!)\),
taken as zero if \(am+di<v\). Thus \(Q\) is integral, its top
coefficient is \(d^m\), and the coefficients of degrees \(r+1\)
through \(bm\) vanish.

For \(\omega^d=1\) put \(E_\omega(z)=e^{-\omega z}Q(\omega z)\).
Let \(T=e^{-z}Q-B\), which vanishes to order at least \(bm+1\).
Since \(B(\omega z)=B(z)\), we have \(E_\omega-B=T(\omega z)\).
Form the Wronskian of \(B\) and the \(d\) functions \(E_\omega\)
with derivative orders \(0,\ldots,d\). Factor \(e^{-\omega z}\)
out of each \(E_\omega\) column. The resulting polynomial
determinant has degree at most \(dm+dr=dbm\).
Its top coefficient comes from row zero of the \(B\) column
and the determinant \(((-\omega)^i)_{1\le i\le d,\,\omega^d=1}\),
which is a Vandermonde determinant times
\(\prod_{\omega^d=1}(-\omega)\). The top coefficient is nonzero.

To calculate the vanishing order, subtract the \(B\) column
from the original \(E_\omega\) columns. The invertible
change of columns taking the sums
\(d^{-1}\sum_{\omega^d=1}\omega^{-(bm+j)}T(\omega z)\),
for \(1\le j\le d\), retains one residue class of exponents in
each column. Their orders are at least \(bm+j\).
In each determinant term the derivatives subtract at most
\(\sum_{i=0}^d i\) from the total order, so the Wronskian vanishes
to order at least \(dbm\). Removing the exponential factors
preserves this order. Hence the polynomial determinant is a
nonzero multiple of \(z^{dbm}\). At any nonzero
root of \(B\) the zeroth row consequently has
\(Q(\omega z)\ne0\) for at least one \(\omega\).
Choose \(G\) to be the minimal polynomial of such a rotation
above a root of \(F\). Then \(\Res(G,Q)\) is a nonzero integer,
and every root of \(F\) has exactly \(k=D/h\) lifts in \(G\).
The norm identity follows by multiplying \(z^d=-x\) over them.

For \(\ell\ge1\), a beta integral gives the tail coefficient
\[
 [z^{bm+\ell}]e^z B(z)
 =\frac{(-1)^m}{m!(\ell-1)!}
       \int_0^1 t^{am}(1-t^d)^m(1-t)^{\ell-1}\,dt.
\]
With \(c_v=[t^v](1+t+\cdots+t^{d-1})^m\), the same coefficient is
\begin{equation}\label{eq:pade-beta}
 (-1)^m\binom{m+\ell-1}{m}
       \sum_v c_v \frac{(am+v)!}{(am+v+m+\ell)!}.
\end{equation}
Let \(p\) be prime and \(e=v_p(z)>0\) in an embedding into
\(\overline{\Q}_p\), normalized by \(v_p(p)=1\), and set
\(\mu=1/(p-1)\). The base-\(p\) digit sum \(s_p(n)\) satisfies
\(s_p(n)/(p-1)\le\log_p(n+1)\), since
\(p^{u/(p-1)}\le u+1\) for each digit \(u\), and the product
of the \(u+1\) counts integers whose digits do not exceed
those of \(n\). If \(e\le\mu\), the coefficient of \(z^{r-j}\) in
\eqref{eq:pade-coeff} contains the integer factor \((m+j)!/m!\).
By Legendre's formula its valuation is at least
\(j\mu-\log_p(M+1)\), since \(m+j\le bm\).
The valuation after multiplication by \(z^{r-j}\) is at least
\(re-\log_p(M+1)\).
If \(e>\mu\), the \(p\)-adic exponential converges and \(Q(z)\)
is the negative tail at a root of \(G\). In \eqref{eq:pade-beta}
put \(n=am+v\le(b-1)m\). Legendre's formula gives
\[
 v_p\!\left(\frac{n!}{(n+m+\ell)!}\right)
 \ge -(m+\ell)\mu-\mu s_p(n)
 \ge -(m+\ell)\mu-\log_p(M+1).
\]
After multiplication by \(z^{bm+\ell}\), the bound increases
with \(\ell\) and has its minimum at \(\ell=1\). In both cases,
\begin{equation}\label{eq:pade-local}
 v_p(Q(z))\ \ge\ (bm+1)e-N\min(e,\mu)-\log_p(M+1).
\end{equation}
For \(e=0\), integrality gives \(v_p(Q(z))\ge0\).

Now let \(z\) be a complex root and set \(R=|z|\).
Integrate each \(v!\) in \eqref{eq:pade-coeff} and use
the nonnegative coefficients of the polynomial in \(w\).
Also sum the beta tail before taking its absolute value.
These two operations give
\begin{equation}\label{eq:pade-integrals}
 \frac{|Q(z)|}{R^{bm+1}}\le\frac1{m!}\min\left\{
 \begin{aligned}
 &\int_0^\infty e^{-Ru}(1+u)^{am}((1+u)^d-1)^m\,du,\\
 &\int_0^1 e^{u\Re z}(1-u)^{am}(1-(1-u)^d)^m\,du
 \end{aligned}\right\}.
\end{equation}
In the second integral \(1-(1-u)^d\le du\). When \(am>\Re z\),
integrating \(u^m e^{-(am-\Re z)u}\) on the positive ray gives
the second gain in the lemma. When \(\Re z\le am\), the integral
is at most one. If \(R\le m\rho_0\), maximizing the exponent
\(u\rho_0+a\log(1-u)\) gives the upper bound
\(e^{m(\rho_0-a-a\log(\rho_0/a))}\). If \(R\ge m\rho_0\),
bound the first integrand using \(((1+u)^d-1)^m\le(1+u)^{dm}\)
and reserve \(e^{-u}\) for integration. The maximum of the
remaining exponent divided by \(m\) is
\(g_+(\rho_0-1/m)\), where \(g_+(v)=b\log(b/v)-b+v\).
Since \(\rho_0-b/2=a/2-s/14>19a/56>1/m\), the interval
\([\rho_0-1/m,\rho_0]\) lies in \((b/2,b)\), where
\(g_+'(v)=1-b/v>-1\). Therefore
\(g_+(\rho_0-1/m)<g_+(\rho_0)+1/m\) and
\(|Q(z)|<R^{bm+1}/L(z)^N\).
Stirling's lower bound gives the simpler gain once
\(0.01(m+1)+\tfrac12\log(2\pi/m)>0\).
This expression increases for \(m\ge199\), since its derivative
is \(0.01-1/(2m)\). At \(199\) it is positive by \(2\pi>6\)
and \(e^4>199/6\), using the first five terms of the exponential series.
When \(d=1\), use the first integral directly, with
\(((1+u)^d-1)^m=u^m\). The tangent inequality
\[
 \log(1+u)\le \log(1+v)-\frac{v}{1+v}+\frac{u}{1+v},
 \qquad v=a^{-1/2},
\]
bounds the constant exponent by \(m/2\), since
\(\log(1+v)-v/(1+v)\le v^2/2\). The coefficient of
\(-u\) is \(R-am/(1+v)>m\sqrt a/(1+v)\). Integrating
\(u^m\) against that exponential proves the additional gain.

Take the product of the complex bounds over the \(D\) roots
and compare with the prime factorization of \(\Res(G,Q)\) using
\eqref{eq:pade-local}. At a prime dividing \(K\) the sum of the
positive \(e\)'s is \(v_p(K)\), and the sum of \(\min(e,\mu)\)
is at most \(\min(v_p(K),D\mu)\). Cancellation of the common
term \((bm+1)\log K\) gives \eqref{eq:pade-norm}.
\end{proof}

For later estimates, the two losses defining \(g\), divided
by \(s\), are the functions
\begin{equation}\label{eq:pade-loss}
 K_+^2\int_0^1\frac{u\,du}{1-\alpha K_+u},
 \qquad
 K_-^2\int_0^1\frac{u\,du}{1-\alpha+\alpha K_-u}.
\end{equation}
The first integrand increases since \(K_+\) increases and
\(1-\alpha K_+u\) decreases. For the second, put
\(D_-=1-\alpha+\alpha K_-u\). Its derivative with respect to
\(\alpha\) is
\[
 \frac{uK_-}{D_-^2}
 \left(K_-(1-u/2)-\frac{1-\alpha}{7}\right),
\]
which is nonnegative for \(0\le u\le1\), since
\(K_-/2\ge25/112>1/7\) on \(0\le\alpha\le3/4\).
At \(\alpha=.03,.1,1/3,.5,.75\), upper bounds for both
functions are respectively \(.128,.134,.160,.176,7/32\).

We also need the following refinement of the local estimate.
Partial fractions in the tail formula yield
\[
 [z^{bm+\ell}]e^z B(z)=
 \frac{(-1)^m d^m}{(\ell-1)!}
 \sum_{j=1}^\ell
 \frac{(-1)^{j-1}\binom{\ell-1}{j-1}}
      {\prod_{i=0}^m(am+di+j)} .
\]
Indeed the coefficient of \(1/(x+j)\) in \(1/(x+1)_\ell\) is
\[
 \left(\prod_{\substack{1\le u\le\ell\\u\ne j}}(u-j)\right)^{-1}
 =\frac{(-1)^{j-1}}{(\ell-1)!}\binom{\ell-1}{j-1}.
\]
Apply this to \(x=am+di\) in the factorial expansion of the
tail. The remaining sum over \(i\) is
\(\sum_i(-1)^i\binom mi/(x+di)=d^m m!/\prod_{i=0}^m(x+di)\).
The coefficient of \(1/(x+di)\) on the right is
\(m!/\prod_{t\ne i}(t-i)=(-1)^i\binom mi\).
In both identities, matching the simple poles leaves a polynomial
that tends to zero at infinity.

For \(e>\mu\), the tail term indexed by \(\ell,j\) has
valuation at least
\[
 (bm+1)e+m v_p(d)
  -v_p\!\left(\prod_{i=0}^m(am+di+j)\right)+(\ell-1)(e-\mu),
\]
using \(v_p((\ell-1)!)\le(\ell-1)\mu\).
At fixed \(j\) the weakest bound is at \(\ell=j\).
Taking the infimum gives
\begin{equation}\label{eq:pade-progression}
 v_p(Q(z))\ge (bm+1)e+m v_p(d)-
 \sup_{j\ge1}\left\{v_p\!\left(\prod_{i=0}^m(am+di+j)\right)
                        -(j-1)(e-\mu)\right\}.
\end{equation}

\subsection{Rotations and the prime support}

Write \(U_i(j)\) for the reverse Newton ordinates
\[
 U_i(j)=v_p((bi-dj+1)_{dj})-v_p\binom ij,\qquad i=N-1,N.
\]
For \(L_p=\lfloor\log_p M\rfloor\), the \(dj\) factorial
arguments contain at most \(dj/p^\nu+1\) multiples of \(p^\nu\)
for \(1\le\nu\le L_p\), and none for \(\nu>L_p\).
Summing and dropping the binomial valuation gives
\(U_i(j)\le dj/(p-1)+L_p\). The concave hull has the
same affine upper bound, so its value at \(h\) bounds the
sum of the largest \(h\) valuations of roots of \(R_i\).
In particular
\begin{equation}\label{eq:norm-affine}
 v_p(F(0))\le H/(p-1)+L_p,\qquad
 p>2H\ \Longrightarrow\ v_p(F(0))\le L_p.
\end{equation}
Also every prime of either initial block \((ai+1)_H\)
divides \(F(0)\). Otherwise \(F\bmod p\), of degree \(h\) and
nonzero constant term, would divide a power of \(x\) times
a polynomial of degree less than \(h\).

\begin{lemma}[Cyclotomic restriction]\label{lem:kummer}
Above any one root of \(F\), the roots of \(G\) in
Lemma~\ref{lem:resultant} are a union of regular rotation cosets
of common size \(k_0\mid d\). If \(q=d/k_0\ge2\), and the
starting-point inequality of Lemma~\ref{lem:degree-cap} holds,
then
\begin{equation}\label{eq:kummer-support}
 2H<P=M^{1/(q+1)},\qquad p\mid F(0)\ \Longrightarrow\ p\le P.
\end{equation}
If \(k_0=d\), then \(k=d\).
\end{lemma}

\begin{proof}
Let \(x\) be a root of \(F\) and adjoin \(\mu_d\) to \(\Q(x)\).
Over this field the extension generated by \(z^d=-x\) is Galois,
and \(\sigma\mapsto\sigma(z)/z\) identifies the group with a
subgroup of \(\mu_d\), of order \(k_0\). Any other lift differs
from \(z\) by a root of unity already adjoined. This proves
the coset assertion and shows that \(z^{k_0}\) is a
\(q\)-th root of \(-x\) in \(\Q(x,\mu_d)\).
At \(p>d\) the cyclotomic extension is unramified, by lifting
roots of unity of order prime to \(p\) from the residue field.
Thus every positive integer-normalized local valuation of \(x\)
in \(\Q(x)\) is a multiple of \(q\). With normalization
\(v_p(p)=1\) it is at least \(q/h\), and \(v_p(F(0))\) is
also a multiple of \(q\).

Suppose \(p>H+d\), \(p^{q+1}>M\), and \(p\mid F(0)\).
For a selected slope \(\lambda\ge q/h\), the affine
upper bound puts every maximizing index of \(U_i(j)-\lambda j\)
below \(qh/(q-1)\le2h\). In this range \(j! \) is a unit.
The factorial jumps are spaced by at least \(h+1\) indices,
with heights at most \(q\), and the binomial valuation cannot
decrease. Thus the maximizing indices can only be zero and
the first factorial jump. By \eqref{eq:norm-affine}, the norm
exponent is exactly \(q\). The edge at each order must have
height \(q\), since it accounts for all selected positive slopes.
The equal widths then place the difference of the two reached
multiples of \(p^q\) in \([a+1,a+2d-1]\). Consequently every norm
prime is at most
\[
 \max\{H+d,\ M^{1/(q+1)},\ (a+2d-1)^{1/q}\}.
\]
The starting-point inequality supplies an initial-block prime
above \(2H\) by Lemma~\ref{lem:degree-cap}. If
\((a+2d-1)^{1/q}>2H\), then
\(s<(1+\alpha)/(4h^2)\), and \(b<dN\) gives
\[
 \frac{M}{(a+2d-1)^{3/2}}>
 \frac{1}{\sqrt s(1+\alpha)^{3/2}}>
 \frac{2h}{(1+\alpha)^2}>1.
\]
It follows that \(M^{1/(q+1)}>(a+2d-1)^{1/q}\) in this case.
If the last entry in the maximum is at most \(2H\), the
initial-block prime directly gives \(M^{1/(q+1)}>2H\).
Both cases prove \eqref{eq:kummer-support}.
\end{proof}

A coset of size \(n\ge2\) places at least \(\lfloor n/2\rfloor\)
points in the closed left half-plane, where the gain is at
least \(m(1-\alpha)/\alpha\). We will use a stronger average
when \(\alpha\) is small.

\begin{lemma}[Angular average]\label{lem:angles}
Assume \(N\ge200\) and \(\alpha<.03\), and put
\(g_0=1.01+.128s\) and \(C'=(1-\alpha)/(2\alpha)\).
In a coset of size \(n\ge2\), order the arguments
\(\theta_j\) by absolute value, and put \(u_j=|1-e^{i\theta_j}|^2\).
Let \(r\) be the number of \(j>0\) with \(u_j<4\alpha\).
The logarithm of the geometric mean of the gains is at least
\begin{equation}\label{eq:pade-angles}
 \log m+(1-1/n)\log C'
 +\frac{2\log(2n/\pi)-g_0}{n}
 -\frac{r}{n}(g_0+\log(2(1-\alpha)))-J,
\end{equation}
where
\[
 \frac rn\le \frac2\pi\arcsin\sqrt\alpha<.111,\qquad
 J=\frac{2\log2}{\pi}
        \bigl(\arcsin\sqrt\alpha+\sqrt{\alpha(1-\alpha)}\bigr)<.153.
\]
For \(n\le9\) one has \(r=0\) and may use
\(J=2(\log2)\alpha(n^2-1)/(3n)\). A weaker bound, used below,
takes \(C=C'/2\) in \eqref{eq:pade-angles}, replaces
\(\log(2(1-\alpha))\) by \(\log(1-\alpha)\), and omits \(J\).
\end{lemma}

\begin{proof}
Where \(u_j\ge4\alpha\), the gain is at least
\(mC'u_j(1-2\alpha/u_j)\). For nonnegative cosine use
\(|z|\le bm\) in \(am-\Re z\). For negative cosine use
\(|z|>am\). Use the baseline \(me^{-g_0}\) at the nearest
argument and at the other \(r\) arguments. The cyclotomic
product gives \(\prod_{j>0}u_j\ge(2n/\pi)^2\), because
\(|\theta_0|\le\pi/n\) and
\(|\sin(n\theta_0/2)/\sin(\theta_0/2)|\ge2n/\pi\).
The removed factors \(C'u_j\) are less than \(2(1-\alpha)\).
The angle spacing gives \(|\theta_j|\ge\pi j/n\) and the
bound on \(r/n\). For the retained factors, apply
\(\log(1-y)\ge-2(\log 2)y\) for \(0\le y\le1/2\).
The mean correction is bounded by the integral of the
decreasing function
\(2(\log2)\min\{1/2,\alpha/(2\sin^2(\pi x/2))\}\) over
\(0\le x\le1\), which equals \(J\).
If \(n\le9\), then \(\sin^2(\pi/(2n))>.03\) and the
discrete sum gives the claimed improvement using
\(\sum_{j=1}^{n-1}\csc^2(\pi j/(2n))=2(n^2-1)/3\).
Finally the difference between the strong and weak bounds
is \((1-(r+1)/n)\log2-J>0\).
\end{proof}

\subsection{Large primes in the norm}

\begin{lemma}[Two possible edge heights]\label{lem:two-heights}
Assume \(d\ge2\) and let \(p>2H+d\), \(p\mid F(0)\), and
\(p^4>M\). Suppose neither
interval \([bi-6H,bi]\), \(i=N-1,N\), contains a positive
multiple of \(p^3\). In particular this holds if \(p^3>M\).
The selected Newton edge at order \(i\) joins zero to the first
jump, of width \(w_i\) and height \(e_i\in\{1,2\}\).
Write \(u_i=bi-dw_i+q_i\), \(1\le q_i\le d\), for the reached
multiple of \(p\). The heights and widths satisfy
\[
 \frac{e_N}{w_N}=\frac{e_{N-1}}{w_{N-1}},\qquad
 v_p(F(0))\le\min(e_N,e_{N-1}).
\]
If the heights are equal, \(u_N-u_{N-1}\in[a+1,a+2d-1]\).
Otherwise \(v_p(F(0))=1\), the widths are in ratio two,
and \(p\) divides an integer in one of
\begin{equation}\label{eq:unequal-products}
 b(N-2)+[2-d,2d-1],\qquad b(N+1)+[2-d,2d-1].
\end{equation}
In this case \(p^2\) has a positive multiple in one of
\([bi-2H,bi]\). In both cases \(u_N\ne u_{N-1}\) and
\[
 u_N-u_{N-1}\in[b-H-d+1,b+H+d-1],
 \qquad p^{v_p(F(0))}\mid (u_N-u_{N-1}).
\]
\end{lemma}

\begin{proof}
The slope \(\lambda\) is at least \(1/h\). The affine
count with error three puts all maximizing indices below
\(6h\), and the absence of terminal cubes improves this to
\(4h<p\). The binomial valuation is nondecreasing there.
A factorial jump raises the ordinate by at most two, and
successive jumps are more than \(2h\) apart. Thus the
weighted ordinates \(U_i(j)-\lambda j\) strictly decrease
between jumps and from one jump index to the next. Only zero
and the first jump can maximize. This proves the first assertions.

The statements for equal heights follow by subtracting the
\(u_i\). For unequal heights the smaller width is the reduced
denominator of the slope and is at most \(h\). Twice the
reached multiple with the shorter width minus the other
belongs to \eqref{eq:unequal-products}. The wider edge
reaches a multiple of \(p^2\) within \(2H\) of the terminal.
In either case, \(|w_N-w_{N-1}|\le1/\lambda\le h\)
gives the stated interval for the difference. If the two
multiples were equal then \(w_N>w_{N-1}\), since \(b>d\).
The upper binomial numerator interval would contain the
lower one, with the same factorial valuation at the first
jump. This gives \(e_N\le e_{N-1}\), contradicting the common
positive slope. The norm bound on \(v_p(F(0))\) gives the last
divisibility.
\end{proof}

\begin{lemma}[Prime sums used below]\label{lem:prime-estimates}
Let \(S(x)=\sum_{p\le x}\log p/(p-1)\). Then for \(x>1\),
\[
 \pi(x)\le\frac{x}{\log x}\left(1+\frac{1.2762}{\log x}\right),
 \qquad S(x)<\log x+.2.
\]
For \(x\ge350\) the latter bound improves to \(\log x-.49\),
and for \(\log x\ge8\) to \(\log x-.54\).
\end{lemma}

\begin{proof}
The prime count is Dusart's Theorem 6.9 \cite{dusart-primes}.
Theorem 6.11 gives an upper bound for
\(\sum_{p\le x}(\log p)/p\) of
\(\log x+E+.2/\log x+.2/\log^2 x\) when \(x>2974\).
Add the prime-power correction using
\(E=-\gamma_{\rm E}-\sum_p\log p/(p(p-1))\) and
\(\gamma_{\rm E}>.57\). This proves the claimed bound
at \(\log x\ge8\). For \(1<x<3000\), the excess \(S(x)-\log x\)
decreases between primes. Enumerate primes by trial division
and bound the excess at each prime using the expansion
\(\log r=2\sum_{j=0}^9 z^{2j+1}/(2j+1)+\epsilon\), where
\(r\in[1,2]\), \(z=(r-1)/(r+1)\), and
\(0\le\epsilon\le2z^{21}/(21(1-z^2))\).
Normalize arguments by powers of two. This gives an upper
bound below \(.144\) on the excess, with its maximum at \(p=3\).
Starting at \(x=350\) and checking that endpoint as well gives
an upper bound below \(-.499\) through \(3000\).
\end{proof}



\section{Completion of the proof}

The support and initial-order reductions leave the range of
Lemma~\ref{lem:support-reduction}. Proposition~\ref{prop:finite}
covers \(25\le N<3150\) there. Proposition~\ref{prop:middle-tail}
covers \(3150\le N<e^{15}\), and
Proposition~\ref{prop:large-tail} covers \(N\ge e^{15}\).
Together with Proposition~\ref{prop:small-terminals}, these give
rational coprimality in every case of \eqref{eq:moment} and hence
Theorem~\ref{thm:main}.

\paragraph{Electronic supplement.}
The electronic supplement accompanying this article contains two parts.
Its exact verification programs cover the initial and central orders,
the finite boundary cases, and the bounded comparisons in the
Pad\'e estimates. The submitted proof appendix
\textit{Positive-exponent Boundary} gives the three contour
arguments used in Appendix~\ref{app:positive-boundary}, with their
exact constants and verification programs.

\paragraph{AI assistance and review.}
ChatGPT (OpenAI) assisted with mathematical arguments, literature checks,
exposition, proof review and verification code.

\appendix
\section{Proofs of the root bounds and additional boundary cases}
\label{sec:gamma-annulus-proof}

\subsection{Bounds for a Gamma polynomial}

For \(\gamma>0\) put
\[
 G_{m,\gamma}(z)=\sum_{k=0}^m\binom{m}{k}(\gamma)_kz^{m-k},
 \qquad h_0=m-1 .
\]
We first show that every zero \(z\) of \(G_{m,\gamma}\) satisfies
\[
 |z|\le\gamma+h_0,\qquad
 \Re z\le-\gamma+h_0.
\]
Let \(T=\gamma+m\). The polynomial is \(\det(zI+J)\), where
\(J=D+iK\) is complex symmetric, with real diagonal
\(D_{jj}=T+1-2j\) for \(1\le j\le m\) and real symmetric
off-diagonal entries \(K_{j,j+1}=\sqrt{j(T-j)}\).
To check the representation, the characteristic polynomial of the
leading \(j\times j\) block is
\[
 \sum_{k=0}^j(-1)^{j-k}\frac{j!}{k!}\binom{T-1-k}{j-k}z^k .
\]
This expression has the initial terms \(1,z-(T-1)\) and
satisfies the determinant recurrence, whose two coefficients are
\(z-(T+1-2j)\) and \((j-1)(T+1-j)\). At \(j=m\), replacing \(z\)
by \(-z\) and multiplying by \((-1)^m\) gives \(G_{m,\gamma}\).

The lower bound on the eigenvalue real parts of \(J\) follows
from its real diagonal \(D\). To bound its numerical radius,
take \(0<\theta<\pi\) and set \(q=\tan(\theta/2)\).
For the Hermitian part \(H_\theta=(\cos\theta)D+(\sin\theta)K\) of
\(e^{-i\theta}J\), the positive vector with entries
\[
 v_j=\sqrt{\binom{T-1}{j-1}}\;q^{j-1}
\]
satisfies \((H_\theta v)_j=(T-1)v_j\) for \(j<m\), and
\((H_\theta v)_m=[T-1-\gamma(1-\cos\theta)]v_m\).
These identities follow by substituting the adjacent ratios
in the off-diagonal terms. The top eigenvector of \(H_\theta\)
can be taken positive by maximizing its Rayleigh quotient
with nonnegative coordinates and using the connected off-diagonal
chain. Its scalar product with the displayed inequality gives
\(\lambda_{\max}(H_\theta)<T-1\). An alternating diagonal
conjugates \(K\) to \(-K\), and continuity handles the endpoint
angles. Hence \(|z|\le T-1\) for each eigenvalue \(z\) of
\(-J\). For \(m\ge2\), equality could only occur in an endpoint
direction. An eigenvector would then be supported in a single
extreme diagonal coordinate, which the adjacent entry prevents.
Thus the modulus bound is strict for \(m\ge2\).

Coefficient comparison also gives the equation
\[
 zG''+(1-\gamma-m-z)G'+mG=0 .
\]
The constant term is positive. At a repeated nonzero zero the
term \(zG''\) would have strictly smaller order than every
other term, so the zeros are simple. Let \(z_i\) minimize
the modulus and write \(r=|z_i|\). The real part of the
differential equation at \(z_i\), using
\(G''(z_i)/G'(z_i)=2\sum_{j\ne i}(z_i-z_j)^{-1}\), gives
\begin{equation}\label{eq:gamma-radius-sum}
 r-\gamma\ \ge\ -\Re z_i-\gamma
  =\sum_{j\ne i}\frac{|z_j|^2-r^2}{|z_j-z_i|^2}\ \ge\ 0 .
\end{equation}
Equality \(r=\gamma\) forces \(z_i=-\gamma\) and all moduli
\(\gamma\). Since their sum is \(-m\gamma\), all zeros would
then equal \(-\gamma\). Thus \(r>\gamma\) for \(m\ge2\).

When \(m\ge2\) and \(\gamma\ge15h_0\), we show the stronger
\(r>\gamma+h_0/6\). Vieta's formulas give
\(\sum z_j=-m\gamma\) and
\(\sum z_j^2=m\gamma^2-mh_0\gamma\). In particular,
\[
 \sum_j |z_j|^2\ge m\gamma^2+mh_0\gamma .
\]
Set \(u_j=z_j+\gamma\) and \(s_0=h_0/\gamma\).
The two bounds on \(z_j\) above give \(|\Re u_j|\le h_0\).
Since \(\sum u_j=0\) and \(\sum u_j^2=-mh_0\gamma\), it
follows that \(\sum |u_j|^2\le mh_0\gamma(1+2s_0)\).
Suppose \(r\le\gamma+h_0/6\), and set
\(\eta=(r^2-\gamma^2)/(\gamma h_0)\). Then
\(0\le\eta\le181/540\) and
\(|u_i|^2\le r^2-\gamma^2\), since \(\Re z_i\le-\gamma\).
Write \(q_j=|z_j|^2-r^2\). We have
\[
 \sum_j q_j\ge m\gamma h_0(1-\eta),\qquad
 0\le q_j\le\gamma h_0(2+s_0-\eta),\qquad
 \sum_j |z_j-z_i|^2\le m\gamma h_0(1+2s_0+\eta).
\]
The last inequality follows from \(\sum u_j=0\).
Weighted Cauchy--Schwarz in \eqref{eq:gamma-radius-sum} now gives
\[
 r-\gamma\ge
 m\frac{(1-\eta)^2}{(2+s_0-\eta)(1+2s_0+\eta)}
 \ge m\,\frac{128881}{741455}>\frac{h_0}{6}.
\]
The rational expression decreases in \(s_0\) and \(\eta\)
on the stated ranges, so its value at \(s_0=1/15,\eta=181/540\)
gives the middle bound. This contradiction proves the stronger
inner radius.

\subsection{Products of Gamma annuli}

For monic polynomials of degree \(m\) in the binomial basis, write
\[
 \left(\sum_{k=0}^m\binom{m}{k}a_k z^{m-k}\right)\diamond
 \left(\sum_{k=0}^m\binom{m}{k}b_k z^{m-k}\right)
 =\sum_{k=0}^m\binom{m}{k}a_kb_k z^{m-k}.
\]
The upper and lower root radii multiply under this operation.
For the upper bound, if the first polynomial is
\(\prod_{j=1}^m(z+\delta_j)\) and \(Q^{\mathrm{pol}}\) is the
polarization of the second, the composed value at \(z\) equals
\(\prod_j\delta_j\, Q^{\mathrm{pol}}(z/\delta_1,\ldots,z/\delta_m)\).
If \(|z|\) exceeds the product of the two upper radii, all the
arguments lie in the open exterior of the disk for the second
polynomial. The Grace--Walsh--Szeg\H{o} theorem
\cite[Theorem~4.1]{wagner-stability} gives nonvanishing there.
The lower bound follows by reversal, which commutes with
\(\diamond\). With strict radii one can first tighten the
bounds on the finitely many roots of each factor.

Use \(\gamma=am+1\) and the identity
\[
 (\gamma)_{dk}=d^{dk}\prod_{j=0}^{d-1}((\gamma+j)/d)_k .
\]
It expresses \(R_m\), after scaling the variable by \(d^d\),
as the composition of the \(d\) Gamma polynomials of shapes
\((\gamma+j)/d\). Consequently
\[
 (\gamma)_d\le |x|\le(\gamma+d(m-1))_d
 \qquad(R_m(x)=0),
\]
strictly for \(m\ge2\). This gives the first assertion of
Lemma~\ref{lem:gamma-annulus}. At the two orders
\(\gamma=am+1\), the inner radius for \(m=N\) is at least
the outer radius for \(m=N-1\) when \(a\ge d(N-2)\).
The strict inner bound gives separation even at equality and
at \(N=2\). If \(N\ge25\) and \(6a\ge(5N-11)d\), then
\(a\ge15d\). Each of the \(d\) unit-step shapes for
\(R_N\) is at least \(15(N-1)\). The stronger bound gives
the lower radius
\((aN+1+d(N-1)/6)_d\), and
\[
 aN+1+d(N-1)/6\ge a(N-1)+1+d(N-2).
\]
This proves the remaining separation in
Lemma~\ref{lem:gamma-annulus}.

\subsection{Small positive lower exponents and adjacent exponents}

\begin{proposition}
\label{prop:additional-boundary}
For \(N\ge25\), the pair \(P_{N-1},P_N\) is coprime on
each of the following supports:
\[
 1\le a\le4,\quad 1\le d<3a,
 \qquad\text{or}\qquad
 a\ge1,\quad d=1.
\]
\end{proposition}

\begin{proof}
Suppose first \(N\ge3150\). We may assume
\(6a<(5N-11)d\), by Lemma~\ref{lem:gamma-annulus}. In both
domains \(d<3a\), and Lemma~\ref{lem:degree-cap} applies
because \(aN\ge1375\). Thus the primes of the common-factor
norm in Lemma~\ref{lem:resultant} are supported up to \(C_N\)
as in Lemma~\ref{lem:prime-support}.

If \(1\le a\le4\) and \(d<3a\), then \(b\le15\),
\(a+2d-1\le25\), and \(s\le121/15\). Use the prime threshold
\(X=\sqrt{15N}\), which is at least \(C_N\) and exceeds \(150\).
By Lemmas~\ref{lem:prime-estimates} and~\ref{lem:small-prime-sums},
the right side of \eqref{eq:pade-norm} is less than
\[
 \tfrac12(\log N+\log 15)-.45+2.48\sqrt{15/N}.
\]
Here \(\log(M+1)/\log X<2.001\), and
\(2.001(1+1.2762/\log X)<2.48\).
The loss bound \(g\le7s/32\) in \eqref{eq:pade-loss} gives
\(\log\Lambda>\log N-2.776\).
The difference exceeds \(0.17\) at \(N=3150\) and increases
with \(N\).

For \(d=1\) and \(a\ge2\) take \(X=\sqrt{(a+1)N}=C_N\),
since \(a+1<(5/6)(N-1)\). The prime-sum and prime-count
bounds give instead
\[
 \tfrac12(\log N+\log(a+1))+.2+2.60\sqrt{(a+1)/N}.
\]
Use the \(d=1\) gain from Lemma~\ref{lem:resultant}.
The surplus is greater than
\[
 \tfrac12\log N-.701-\tfrac12\log(3/2)
       -\log(1+1/\sqrt2)-2.60\sqrt{5/6}>.21
 \qquad(N\ge3150).
\]
The case \(a=d=1\) is in the preceding domain.

For \(25\le N<3150\), the supplemental checker
\texttt{verify\_small\_support.cpp} treats the two stated
support families together. It retains the inequalities
\(a<d(N-2)\), \(6a<(5N-11)d\) and \(d^2<19(a+d)\), using
Lemmas~\ref{lem:gamma-annulus} and~\ref{lem:ratio-range},
and uses the gap caps in Lemmas~\ref{lem:gap-caps}
and~\ref{lem:local-gap}. It groups the first factorial block
test by \((N,a)\) and then applies the individual block and
Newton slope tests of Section~\ref{sec:arithmetic}.
The resulting counts are \(12\,520\) groups and \(18\,904\)
supports. Of these, \(4\,357\) are excluded by the degree-dependent
prime bounds and \(4\,853\) by comparing Newton polygons.
For the other \(9\,694\), the full monic reciprocals are coprime
modulo \(1\,000\,000\,007\). The checker verifies primality by
trial division and computes the modular gcd by exact Euclidean
division. A monic common factor over \(\mathbb Q\) would be
integral and would retain its degree modulo this prime.
Thus these tests establish coprimality throughout the finite range.
\end{proof}


\section{The constant-plus-power boundary}
\label{app:constant-proof}

Write \(F_{n,d}(t)=\sum_{k=0}^n\binom nk(dk)!t^k\) and
\(R_n(z)=z^nF_{n,d}(1/z)\). We prove that \(R_n\) is squarefree
for all positive integers \(n,d\). The proof uses a coefficient
bound for large \(d\), Newton residuals for the small prime powers,
and a uniform contour estimate for the remaining large orders.

\subsection{The coefficient bound}

For integers \(1\le u\le v\) one has
\[
 \frac{u!}{(v!)^{u/v}}\le (u/v)^{u/2}.
\]
Indeed \((k!)^{1/k}/\sqrt{k}\) is nondecreasing. The comparison
at \(k,k+1\), after raising to positive powers, is
\((k!)^2(k+1)^{k(k-1)}\le k^{k(k+1)}\). This follows from
\(k!\le((k+1)/2)^k\) and \((1+1/k)^{k+1}\le4\). For the latter
bound, \((x+1)\log(1+1/x)\) decreases for \(x>0\).
Also \(k\log(n/k)\ge(\log2)\min(k,n-k)\), by concavity of
\(-x\log x\) on \([1/2,1]\) and by \(n/k\ge2\) for \(k\le n/2\).

After scaling \(t\) by \(((dn)!)^{-1/n}\), the polynomial is
\(1+t^n+\sum_{k=1}^{n-1}b_kt^k\), with
\[
 0<b_k\le C_k:=\binom nk q^{\min(k,n-k)},\qquad q=2^{-d/2}.
\]
The bound follows from the factorial inequality at \(dk,dn\).
If \(2^d\ge4n^2\), then \(q\le1/(2n)\).
By symmetry of \(C_k\) and \(q^{\min(k,n-k)}\le q^k+q^{n-k}\),
both weighted sums \(\sum(k/n)b_k\) and
\(\sum(1-k/n)b_k\) are bounded by \((1+q)^n-1<1\).
At a repeated root of modulus at most one, subtract \(t/n\) times
the derivative from the polynomial to obtain a contradiction to
the second bound. At a repeated root of modulus at least one,
divide \(t/n\) times the derivative by \(t^n\) and use the first.
We have proved that
\begin{equation}
 2^d\ge4n^2 \ \Longrightarrow\  F_{n,d}\text{ is squarefree}.
 \label{eq:constant-large-gap}
\end{equation}

\subsection{Small exponents}

For \(d=1\), \(R_n(z)=n!\sum_{j=0}^n z^j/j!\). Its derivative
and itself cannot share a root, since the two truncated
exponentials differ by \(z^n/n!\). To cover \(2\le d\le5\),
we give the prime-power argument. The factorial units and the
digit structure of the Newton polygon used here occur in
Rocks's work on factorial moment polynomials~\cite{rocks}.

Let \(d=p^\nu\) with \(p\) prime and \(\nu\ge1\),
and put \(c=(d-1)/(p-1)\). For an integer \(u\ge0\), write
\(U_p(u!)=p^{-v_p(u!)}u!\). If \(u=\sum_i u_i p^i\),
Wilson's theorem, applied after taking out multiples of \(p\)
one level at a time, gives
\[
 U_p(u!)\equiv(-1)^{v_p(u!)}\prod_i u_i!\pmod p.
\]
The digit shift from \(k\) to \(dk\) and Legendre's formula
therefore give
\[
 F_{n,d}(T/p^c)=\sum_{k=0}^n \frac{n!}{(n-k)!}w_k T^k,\qquad
 w_k=\frac{(dk)!}{p^{ck}k!}\in\Z_p^\times,\qquad
 w_k\equiv(-1)^{ck}\pmod p .
\]
Write \(n=\sum_i a_i p^i\) and \(s_i=\sum_{j\le i}a_jp^j\),
with \(s_{-1}=0\). The lower Newton polygon has an edge from
\(s_{i-1}\) to \(s_i\) for each \(a_i>0\), of slope \(h_i/e_i\)
with \(e_i=p^i\) and \(h_i=(p^i-1)/(p-1)\). Indeed the
coefficient ordinate is
\(v_p(n!)-v_p((n-k)!)\). For \(k=s_i-\delta\) in that segment
the difference between the ordinate and the line through
the two endpoints is
\[
 \frac{\operatorname{digitsum}_p(\delta)-\delta/p^i}{p-1}\ge0.
\]
This is Legendre's formula, since \(n-s_i\) and \(\delta\)
have disjoint digit positions. Equality occurs exactly at the
indices spaced by \(e_i\), the slopes strictly increase,
and \(h_i/e_i\) is in lowest terms.

Divide each coefficient on the edge by its power of \(p\)
and take ratios to the leftmost coefficient. Wilson's unit
formula then gives the edge polynomial
\[
 \sum_{j=0}^{a_i}\frac{a_i!}{(a_i-j)!}(\eta_i Y)^j,\qquad
 \eta_i=(-1)^{h_i+c e_i}.
\]
Here only digit \(a_i\) changes in the factorial ratio, and the
change in valuation is \(jh_i\). The displayed polynomial
is a scaled reversal of \(E_{a_i}(Z)=\sum_{j=0}^{a_i}Z^j/j!\)
over \(\mathbb F_p\). It is squarefree because
\(E_{a_i}'=E_{a_i-1}\) and \(E_{a_i}-E_{a_i-1}=Z^{a_i}/a_i!\).

For completeness, edge polynomials multiply at any fixed
slope \(u/e\) in lowest terms. Substitute \(T=\pi^{-u}Z\)
in a valued extension with \(\pi^e=p\), divide each polynomial
by the power giving minimum coefficient valuation zero,
and reduce. The reductions multiply. The surviving indices
in a polynomial over \(\Q_p\) are all congruent modulo \(e\),
so after removing the lowest monomial the same identity
holds in the edge variable \(Y=Z^e\).
A square divisor would thus give a repeated factor in the
edge polynomial at any one of its slopes. The squarefreeness
above rules this out.

\subsection{A uniform contour estimate}

We now assume \(n\ge10000\) and \(6\le d\le\sqrt n\).
Set \(A=(d/e)^d\) and \(v=z/(n^d A)\).
The exponential integral for \(R_n\), with its contour split
at \(d/2\), is
\[
 R_n(n^d A v)=n^{dn+1}I(v),\qquad
 I(v)=\int_0^\infty(Av+s^d)^n e^{-ns}\,ds=K(v)+J(v),
\]
where \(K\) and \(J\) integrate over \([0,d/2]\) and
\([d/2,\infty)\), respectively. Both are polynomials in \(v\).
We bound \(J\) by a Gaussian at its saddle and \(K\) by an
endpoint approximation. Their logarithmic derivatives are
separated where a zero of \(I\) can occur.

\paragraph{The outer bound.}
For \(\theta=k/n\), the entropy bound on \(\binom nk\) and
\((dk)!\le e\,dk(dk/e)^{dk}\) give
\[
 \frac{[\binom nk(dk)!]^{1/k}}{n^d A}
 \le q_d(\theta)(e\,dk)^{1/k},\qquad
 q_d(\theta)=\theta^{d-1}(1-\theta)^{-(1-\theta)/\theta}.
\]
The endpoint values use continuous limits. We have \(q_d\le q_6\).
The sign of the derivative of \(\log q_6\) is the sign of
\(6\theta+\log(1-\theta)\). Its positive-to-negative change
occurs at \(\theta=1-u\) with \(0<u<1/300\), since
\(6(299/300)>\log300\).
There \(\log q_6=5\log(1-u)+6u<u\), so \(q_6<300/299\).
If \(k\ge n/10\), then \(d\le\sqrt n\) gives
\[
 k^{-1}\log(e\,dk)\le10(1+\tfrac32\log n)/n<2/125.
\]
Thus the normalized coefficient root is less than
\((300/299)(125/123)<21/20\).
If \(k<n/10\), the bounds
\(\binom nk^{1/k}\le en/k\) and \((e\,dk)^{1/k}\le ed\)
give instead \(e^2 d\theta^{d-1}<21/20\).
The latter unit bound uses \(k\le2^{k-1}\le(ed)^{k-1}\).
Summing the resulting finite geometric series excludes
\(|v|\ge R=21/10\).

\paragraph{The large saddle.}
Put \(B=11/5\) and \(\delta=1/25\). For \(|v|\le B\) the
equation \(A v=s^{d-1}(d-s)\) has a unique solution
\(s_0=\alpha_0+i\beta_0\) satisfying \(|s_0-d|<\delta\).
On \(|s-d|=\delta\) the modulus on the right divided by \(A\)
is at least \(\delta(1-\delta)e^d/d>B\).
Bernoulli's inequality gives the lower estimate, and \(e>27/10\)
verifies the strict comparison at \(d=6\), where the bound
is smallest. Rouch\'e's theorem gives the solution, which
is simple and analytic in \(v\) on a neighborhood of the disk.
Use the phase
\[
 \Phi(s,v)=d\log s+\log(1+A v/s^d)-s,\qquad
 S=\Re\Phi(s_0,v),\qquad a_0=1-(d-1)/s_0 .
\]
At \(s_0\), \(\Phi_s=0\), \(\Phi_{ss}=-a_0\), and
\begin{equation}
 \Re a_0>9/(10d),\qquad |a_0|<11/(10d),\qquad
 S>\log A-1/100 .
 \label{eq:constant-saddle-bounds}
\end{equation}
For the first two bounds use
\(|a_0-1/d|=(d-1)|s_0-d|/(d|s_0|)<\delta/d\).
For the height, use
\(\Phi(s_0,v)=(d-1)\log s_0+\log d-s_0\).
The loss is at most \(\delta/d+\delta^2/[2d(1-\delta/6)]<1/100\)
by the logarithm series.

Deform the contour of \(J\) to its horizontal translate
by \(i\beta_0\), with the vertical connector at \(d/2\).
The integrand is entire and the connector at positive
infinity tends to zero. On \(s=s_0+t,\ |t|\le d/8\)
we have \(|s|>6d/7\). With \(q_t=B(35/81)^6\),
differentiating the convergent logarithm series gives
\[
 d^2|\Phi_{sss}|
 \le (7/6)^3\left[2+
 \frac{36q_t(1+q_t)}{(1-q_t)^3}
 +\frac{18q_t}{(1-q_t)^2}+\frac{2q_t}{1-q_t}\right]<8.
\]
For this inequality use
\(|Av/s^d|\le B(35/81)^d\).
In the differentiated series the corresponding terms have
factors \(d^2,d,1\), each giving its largest bound at \(d=6\).
The cubic remainder is at most \(4|t|^3/(3d^2)\).
By \eqref{eq:constant-saddle-bounds} its contribution to the local integral is at most
\[
 \frac{4ne^{nS}}{3d^2}\int_{\R}|t|^3 e^{-17nt^2/(60d)}\,dt
 <17e^{nS}/n.
\]
Completing the Gaussian costs less than
\(18e^{nS-9dn/1280}/n\), by bounding the Gaussian tail by
the integral with the extra factor \(8|t|/d\).
The complete value
\(U_0=e^{n\Phi(s_0,v)}\sqrt{2\pi/(na_0)}\)
has modulus greater than \(2\sqrt{d/n}\,e^{nS}\).

We give the bounds on the remaining pieces of this contour.
For \(2d/3\le x\le4d/3\), let \(q_c=B(5/9)^6\).
On \(s=x+i\beta_0\), the logarithm series and
\(|\beta_0|/x<1/100\) give
\[
 \frac{d^2}{dx^2}\Re\Phi(x+i\beta_0,v)
 \le-\frac d{x^2}\left[
 \frac{1-10^{-4}}{(1+10^{-4})^2}
 -\frac{6q_c}{(1-q_c)^2}-\frac{q_c}{1-q_c}\right]
 <-\frac3{16d}.
\]
Here \(|Av/s^d|\le B(5/9)^d\), and \(d(5/9)^d\) decreases
for \(d\ge6\). Integrating from \(x=\alpha_0\) bounds the
height by \(S-3(x-\alpha_0)^2/(32d)\).
Outside the local interval this costs at least \(d/1000\)
in height, over length less than \(d\).
At \(4d/3\) the loss is greater than \(d/100\).
For \(x\ge4d/3\), the real derivative is at most
\((3/4)/(1-B(5/18)^6)-1<-1/5\), so the far ray costs at
most \(5e^{nS-dn/100}/n\).

On \([d/2,2d/3]\) the upper envelope
\([BA+(x+\delta)^d]e^{-x}\) increases. For
\(x+\delta<d-1\), the function
\((x+\delta)^{d-1}(d-x-\delta)\) increases, and its value
at \(x=d/2\) exceeds \(BA\) since \((e/2)^d>B\).
The envelope divided by \(A\) is at most
\(B e^{-2d/3}+r_\ell^d\) there, with
\(r_\ell=(101/150)(141/100)\) from \(e^{1/3}<141/100\).
The corresponding bound on the vertical connector is
\(B e^{-d/2}+r_v^d\) with \(r_v=38/45\), from \(e^{1/2}<5/3\).
Both bounds are below \(e^{-d/50}\).
After multiplication by \(e^{d/50}\) their two exponential
terms decrease in \(d\), and at \(6\) the comparisons follow
from
\[
 B(10/27)^3+r_\ell^6/(22/25)<1,\qquad
 B(10/27)^2+r_v^6/(22/25)<1.
\]
By \eqref{eq:constant-saddle-bounds} the height of these pieces is below \(S-d/100\).
Their lengths are \(d/6\) and at most \(\delta\).
The phase charts used on the horizontal contour are valid
since \(|Av/s^d|\le B(2/e)^d<1\) for \(x\ge d/2\).

Summing the contour bounds gives an error at most \(e^{nS}\)
times
\[
 40/n+d e^{-dn/1000}+(d+\delta)e^{-dn/100}<44/n .
\]
The last comparison uses \(e^x\ge x^3/6\) in the middle
term and \(e^x\ge x^2/2\) in the last.
Consequently \(J=U_0(1+\epsilon)\) with
\(|\epsilon|<22/\sqrt{dn}<1/10\).
The error is analytic on a neighborhood of the disk.
Cauchy's estimate at radius \(1/10\) gives
\(|\epsilon'|<1\) on \(|v|\le R\). Also
\[
 U_0'/U_0=n\frac{d-s_0}{dv}-\frac{a_0'}{2a_0},\qquad
 s_0'=\frac{A}{d s_0^{d-2}(d-1-s_0)} .
\]
Bernoulli's inequality gives \(|s_0'|\le de^{-d}/(1-\delta)^2<1\),
and hence \(|a_0'/(2a_0)|<1\) by \eqref{eq:constant-saddle-bounds}. We obtain
\begin{equation}
 \left|J'/J-n(d-s_0)/(dv)\right|<3 \qquad (|v|\le R).
 \label{eq:constant-large-derivative}
\end{equation}
The expression at zero is understood by continuity.

\paragraph{The endpoint.}
For \(0\le x\le d/2\),
\[
 \log(1+x^d/((2/5)A))\le19x/20.
\]
At \(x=d/2\) this is \(1+(5/2)(e/2)^d<e^{19d/40}\).
After multiplication by \(e^{-19d/40}\), both terms
decrease with \(d\), since \(e^{21/40}\le(5/3)(40/39)<2\).
At \(d=6\) use \(e<49/18\) and the rational comparison
\[
 1+(5/2)(49/36)^6<\sum_{j=0}^{12}(57/20)^j/j! .
\]
To pass to smaller \(x\), the logarithm divided by \(x\)
has derivative of the sign of \(du/(1+u)-\log(1+u)\),
where \(u=x^d/((2/5)A)\). This function is nonnegative up
to \(u=d-1\) by differentiation. Thereafter its first term
is at least \(d-1>19d/40\), while the logarithm is bounded
by its endpoint value.

For \(|v|\ge2/5\), write \(K=(Av)^n(1+D(v))/n\).
Subtract the full exponential integral and use
\(|(1+z)^n-1|\le n|z|(1+|z|)^{n-1}\).
For \(D'\) differentiate \((1+x^d/(Av))^n\) inside
the integral. The same endpoint envelope gives
\[
 |D(v)|\le \frac{d!20^{d+1}}{(2/5)A n^{d-1}}+e^{-dn/2},\qquad
 |D'(v)|\le \frac{d!20^{d+1}}{(2/5)^2 A n^{d-1}}.
\]
Both are below \(1/1000\), since \(d!\le ed A\), with the
largest majorants at \(d=6,n=10000\). In particular
\begin{equation}
 K(v)\ne0,\qquad |K'/K-n/v|<1\qquad(|v|\ge2/5).
 \label{eq:constant-endpoint-derivative}
\end{equation}

For \(|v|\le1/2\), the upper envelope
\((A/2+x^d)e^{-x}\) attains its maximum on \([0,d/2]\)
at an endpoint, by the sign of its derivative.
At the right endpoint it is below \(A/2\) because
\(e^{-d/2}+2(e^{1/2}/2)^d\le(5/8)^6+2(5/6)^6<1\).
Thus \(|K|\le(d/2)(A/2)^n\).
The lower estimate for \(J\) gives
\(|J|>(3/2)\sqrt{d/n}(A e^{-1/100})^n\).
It follows that \(|K/J|<\sqrt{dn}(50/99)^n/3<1\),
using \(\sqrt{dn}\le n\) and \(\sum_{j\ge1}j(50/99)^j<3\).
Therefore every zero of \(I\) lies in \(1/2<|v|<R\).
At such a zero, \(I'=J(J'/J-K'/K)\), and
\eqref{eq:constant-large-derivative} and
\eqref{eq:constant-endpoint-derivative} give
\[
 |J'/J-K'/K|>n|s_0|/(d|v|)-4>n/3-4>0 .
\]
This proves simplicity in the entire stated contour range.

\subsection{Completion and finite verification}

If \(n\ge10000\) and \eqref{eq:constant-large-gap} fails,
then \(d<2\log_2(2n)<\sqrt n\). The second inequality
follows since \(u-2-4\log_2 u\) is positive at \(u=100\)
and increasing for \(u\ge100\). Hence the contour estimate
and the small-exponent argument cover every such order.
For \(n<10000\), the remaining pairs have \(6\le d\le29\)
and \(2^d<4n^2\). The case \(n=1\) is linear.

Here is a small-degree test for that finite range.
Let \(n=p+r\) with \(p\) prime and \(p>d\max(1,r)\).
Then \(r<p\) and \(p^2-dn=p(p-d)-dr>0\).
The Newton ordinate of \(F_{n,d}\) is
\[
 y(k)=\lfloor dk/p\rfloor+\mathbf1_{\{k\bmod p>r\}}.
\]
It is zero up to \(r\). For \(r<k<p\) it exceeds \(dk/p\),
and for \(p\le k\le n\) it equals \(d\).
Thus the polygon has a horizontal part of width \(r\)
and a slope \(d/p\) of width \(p\).
Every repeated rational irreducible divisor has degree
at most \(n/2<p\), so by the slope denominator it can
use only the horizontal part. It is monic integral in
the reciprocal and reduces with nonzero constant term.
Lucas's congruence gives \(R_n\bmod p=z^p(R_r\bmod p)\).
Indeed for \(k\le r\) the binomial coefficient reduces
to \(\binom rk\), and the other coefficients of \(F_{n,d}\)
vanish modulo \(p\). Thus a repeated divisor would
reduce to a repeated divisor of \(R_r\bmod p\).

Every repeated rational irreducible divisor has degree
at least two. For even \(n\), the exponential integral
\(F_{n,d}(t)=\int_0^\infty(1+tu^d)^ne^{-u}\,du\) is positive
on the real axis. For odd \(n\), its derivative is positive there.
In either case it has no repeated real root.
It follows that a gcd of degree at most one between the modular
\(R_r\) and its derivative excludes every repeated rational
factor of \(R_n\). The same argument applies to a direct
modular test of \(R_n\) at any prime greater than \(dn\).

The supplementary verifier exhausts \(202\,021\) pairs in
the remaining range. It uses a shifted prime in \(201\,927\)
cases, where the tested degree is at most \(35\). In the
other \(94\) cases it tests the full polynomial at
\(1\,000\,000\,007\), with degree at most \(148\). The primes
are obtained by an Eratosthenes sieve in the shifted tests
and by trial division for the full tests. Every modular gcd
has degree at most one. This completes the squarefreeness
proof and thus Lemma~\ref{lem:constant-support}.

\section{The large relative gaps}
\label{app:positive-boundary}

We separate the prime-power gaps from the three contour estimates.
The arithmetic proof for prime-power gaps is given below. The
electronic appendix \textit{Positive-exponent Boundary} gives the
quantitative contour bounds, their derivations and the exact
constant checks.

\begin{lemma}
\label{lem:positive-boundary}
The pair \(P_{N-1,a,d},P_{N,a,d}\) is coprime for every \(N\ge2\)
on \(a=1,d\ge3\), and on \(a\ge2,d\ge3a\).
\end{lemma}

\begin{proof}
Proposition~\ref{prop:small-terminals} covers \(N\le24\).
Thus we may impose the ratio bound of Lemma~\ref{lem:ratio-range}
throughout this proof.

For \(a=1\), write
\[
 Q_{m,\gamma,d}(x)=\sum_{k=0}^m\binom{m}{k}
                        (\gamma)_{dk}x^{m-k}.
\]
At terminal \(N\) the actual shapes are \(N\) and \(N+1\).
Coefficient comparison gives, with \(\gamma=N\),
\[
 Q_{N,\gamma,d}'=NQ_{N-1,\gamma,d},\qquad
 \gamma Q_{N,\gamma+1,d}
  =(\gamma+dN)Q_{N,\gamma,d}-dxQ_{N,\gamma,d}'.
\]
In fact the lower reciprocal is \(Q_{N-1,N,d}\), so at a
common root the derivative of \(Q_{N,N,d}\) vanishes. The second
identity then makes the upper reciprocal a nonzero scalar multiple
of \(Q_{N,N,d}\) there. Conversely, a multiple root of
\(Q_{N,N,d}\) makes both reciprocals vanish. The prime-power
argument in the next subsection gives squarefreeness at
\(d=3,4,5,7,8,9\). The separate
gap-six contour proof gives it at \(d=6\) for \(N\ge10000\).

For \(d\ge10\), the lower-one contour proof gives squarefreeness of
\(Q_{N,N,d}\) whenever \(N\ge\max(5000,128(d+1))\).
Its two contour contributions have separated logarithmic
derivatives at a zero, so the derivative of their sum is nonzero.
For \(a\ge2,d\ge3a\), the weighted contour proof in the electronic
appendix gives coprimality for \(N\ge\max(4000,128(a+d))\).
Here one must compare two integrals with the actual moment orders.
With \(w=A/(BN^d)\) and \(g(x)=x^a(w+x^d)\), their integrands
after the common dilation are \(g(x)^N e^{-Nx}\) and
\(g(x)^{N-1}e^{-Nx}\). Write \(I_U,I_L\) and \(J_U,J_L\)
for their two contour contributions. At a zero of the first
integral, \(I_U=-I_L\ne0\). The verified packet estimates at the
two selected saddles \(x_U,x_L\) give
\[
 \left|\frac{J_U}{I_U}\right|
 <\frac{24\sqrt{N/(N-1)}}{|g(x_U)|},\qquad
 \left|\frac{J_L}{I_L}\right|
 >\frac{\sqrt{N/(N-1)}}{5|g(x_L)|},\qquad
 \left|\frac{g(x_U)}{g(x_L)}\right|>150 .
\]
Thus \(|J_U|<|J_L|\) and the second integral is nonzero.
The proof estimates the extra amplitude \(e^{-x}\) in the
order-\(N-1\) contour integral when the phase is written as
\(\log g(x)-x\).

If \(N<\max(4000,128b)\), where \(b=a+d\le4d/3\), a
potential failure has \(d\le24\). Indeed \(s\ge3d/4\).
For \(N\ge200\) and \(d\ge25\), one has
\[
 N<512d/3,\qquad \frac{3d}{4}>2\log(512d/3)+2 .
\]
At \(d=25\) the margin exceeds \(3/100\), and its derivative
\(3/4-2/d\) is positive. For \(N<200\), use
\(\sigma_N\le 2\log 199+3<14<3d/4\).

For \(a=1,d\ge10\) below the contour threshold,
a potential failure has \(d\le19\).
For \(20\le d\le38\) the threshold is \(5000\), and
\[
 \frac{d^2}{d+1}\ge\frac{400}{21}>2\log(5000)+2 .
\]
For \(d\ge39\) the threshold is \(128(d+1)\). The difference
\(d^2/(d+1)-2\log(128(d+1))-2\) exceeds \(18\) at \(39\),
and its derivative \(1-1/(d+1)^2-2/(d+1)\) is positive.
For \(N<200\) and \(d\ge20\), the upper bound
\(\sigma_N<14<400/21\) suffices. The electronic comparison
script certifies the displayed logarithmic inequalities by
rational bounds.

The supplemental finite checker therefore needs only the families
\[
\begin{array}{c|c|c}
a&d&\text{terminals }N\\ \hline
1&6&25\le N<10000\\
1&10\le d\le19&25\le N<5000\\
2\le a\le\lfloor d/3\rfloor&6\le d\le24&
                25\le N<\max(4000,128(a+d)).
\end{array}
\]
Before any further exclusions this is \(338\,071\) cases.
Every case satisfies \(a<d(N-2)\), and a possible common factor
has \(s<21\) by Lemma~\ref{lem:ratio-range}. The block-prime
support and Newton allocation arguments in
Lemmas~\ref{lem:prime-support} and~\ref{lem:newton-alloc}
apply here with these hypotheses. The stronger degree bound
is used only when its starting-point condition is verified
for that degree. In the subsequent Newton check, each odd
prime \(p>aN\) satisfies \(p>N,p>d\) and \(p^2>(a+d)N\)
throughout these finite families. Thus the sparse floor
ordinates used in Section~\ref{sec:arithmetic} are exact
also in this range.

The supplemental program \texttt{verify\_positive\_boundary.cpp}
retains \(326\,146\) cases after imposing the uniform bound
\(d^2<21(a+d)\). It tries the
first available sparse Newton test, applies the prime-block
bounds at the remaining degrees, and then uses additional
primes as needed. The exact exclusions are
\[
\begin{array}{l|r|r|r}
\text{Family}&\text{Supports}&\text{Prime blocks}&\text{Newton polygons}\\ \hline
a\ge2,\ d\ge3a &266\,421&12\,901&253\,520\\
a=1,\ 10\le d\le19&49\,750&966&48\,784\\
a=1,\ d=6&9\,975&535&9\,440
\end{array}
\]
Every possible even degree is excluded, and no modular fallback
is needed.
\end{proof}

\subsection{The prime-power gaps}

We prove the squarefreeness needed above for \(q=p^r\),
where \(p\) is prime and \(r\ge1\). Use \(n=N\), \(\gamma=n\)
and
\[
 T(t)=t^n Q_{n,\gamma,q}(1/t),\qquad
 R(z)=z^{qn}T((-1)^{q+1}z^{-q}),\qquad M=\gamma+qn.
\]
Both changes of variable are unramified away from zero, and both
\(T(0)\) and \(R(0)\) are nonzero. Write \(D=d/dz\), \(\theta=zD\)
and define
\[
 \mathcal A=z-zD+\gamma+q-1,\qquad
 \Pi_{n,q}(t)=\prod_{j=0}^{n-1}(t+qj),\qquad
 B_n=\Pi_{n,q}(\mathcal A)1 .
\]

\paragraph{The rank condition.}
The polynomial \(R\) solves \(LR=0\) for the monic operator
\[
 L=z^{-q-1}\left[
 \theta\prod_{\ell=0}^{q-1}(\theta-M-\ell)-z^q(\theta-qn)\right].
\]
Indeed consecutive coefficients \(r_k=[z^{q(n-k)}]R\) satisfy
\((k+1)r_{k+1}=(-1)^{q+1}(n-k)(\gamma+qk)_q r_k\),
which is the equation at \(z^{q(n-k)}\) after multiplying by
\(z^{q+1}\). Put \(U=1-D\). We will use the adjoint solutions
\[
 f_\omega(z)=e^{-\omega z}z^{1-M}B_n(\omega z),
 \qquad \omega^q=1.
\]
To derive them, set
\[
 \begin{aligned}
 C(x)&=(1-x)^{-\gamma}
       \left(\frac{1-(1-x)^q}{qx}\right)^{-n-1},\\
 A_0(x)&=(1-x)((1-x)^q-1),\\
 b_0(x)&=(M-1)((1-x)^q-1)+qn.
 \end{aligned}
\]
The residue \(n![x^n](C(x)e^{zx})\) equals \(B_n\).
In fact the substitution \(v=1-(1-x)^q\) makes it
\[
 q^n n![v^n](1-v)^{-(\gamma+q-1)/q}
                  e^{z(1-(1-v)^{1/q})}.
\]
With \(v=qt\) the generating function in \(n\) is
\((1-qt)^{-(\gamma+q-1)/q}e^{z(1-(1-qt)^{1/q})}\).
Writing it as \(F(t,z)\), differentiation gives
\((1-qt)F_t=\mathcal A F\), \(F(0,z)=1\). This proves the
claimed residue formula. Logarithmic differentiation gives
\(A_0 C'=((n+1)A_0/x-A_0'+b_0)C\). Taking the residue of
\((e^{zx}A_0 C x^{-n-1})'\) therefore yields
\[
 [zA_0(D)+b_0(D)]B_n
 =z(U^{q+1}-U)B_n+(M-1)(U^q-1)B_n+qnB_n=0.
\]
The divergence form \(L=z^{-q}D z^{M+q}D^q z^{-M}-D+qn/z\)
gives \(L^*=(-1)^{q+1}z^{-M}D^q z^{M+q}D z^{-q}+D+qn/z\).
For \(f_1=e^{-z}z^{1-M}B_n\) the first term becomes
\(e^{-z}z^{-M}[zU^{q+1}+(M-1)U^q]B_n\), using
\(U^qz=zU^q-qU^{q-1}\). The last two terms give
\(e^{-z}z^{-M}[-zU+qn+1-M]B_n\).
This proves \(L^*f_1=0\), and substitution of \(\omega z\)
gives the remaining solutions.

Write \(L=\sum_{k=0}^{q+1}a_kD^k\). The bilinear expression
\[
 \sum_{k=1}^{q+1}\sum_{i=0}^{k-1}
 (-1)^i(a_k f_\omega)^{(i)} R^{(k-1-i)}
\]
has derivative \(f_\omega LR-RL^*f_\omega=0\), by cancellation of
adjacent summands. It is zero because it is a rational function of
\(z\) times \(e^{-\omega z}\), which cannot be a nonzero constant.
At a multiple root \(z_0\ne0\) of \(R\), the first two coordinates
of its jet vanish, and its coordinates \(R'',\ldots,R^{(q)}\)
are not all zero by uniqueness for the ordinary equation \(LR=0\).
They give a common nonzero nullvector of the restricted rows of
these expressions. In the column multiplying \(R^{(j)}\), \(2\le
j\le q\), the term of highest derivative of \(f_\omega\) is
\((-1)^{q-j}f_\omega^{(q-j)}\), since \(a_{q+1}=1\).
The restricted rows are therefore triangular transforms of the
jets \(f_\omega,\ldots,f_\omega^{(q-2)}\). Removing the common
power gauge and the nonzero exponential row factors gives a column
dependence among \(\omega^j(U^jB_n)(\omega z_0)\), \(0\le j\le q-2\).
The relation \(Uz=zU-1\) shows that \((zU)^j\) differs from
\(z^jU^j\) by a combination of \(z^iU^i\), \(i<j\).
Using \(\mathcal A=zU+\gamma+q-1\) and \(z_0\ne0\), the dependence
becomes
\[
 [\Pi_{n,q}(\mathcal A)P(\mathcal A)1](\omega z_0)=0
 \quad\text{for every }\omega^q=1
\]
for some nonzero \(P\) of degree at most \(q-2\).
The dependent evaluation matrix has algebraic entries, so this
condition transfers to \(\overline{\Q}_p\).

\paragraph{Excluding the dependence.}
Normalize \(v(p)=1\), work in a complete algebraically closed
extension of \(\Q_p\), and let \(\mathcal O\) be its valuation ring.
Set \(b_r=1+p+\cdots+p^{r-1}\). Every interval of \(qk\) consecutive
integers contains \(qk/p^\ell\) multiples of \(p^\ell\) for
\(1\le\ell\le r\). Thus \(v((\gamma)_{qk})\ge b_r k\), and
\(v(z_0)\ge b_r/q\). Otherwise the leading term of \(R(z_0)\)
would have uniquely least valuation. In particular
\(v(u)\ge b_r\) for \(u=z_0^q\).

Here is the invariant under descent from \(q=p^r\) to \(p^{r-1}\).
Allow an operator and seed
\[
 A=c-zD+z+\sum_{i\ge1}a_i z^{p^i},\quad
 c\in\Z_p,\quad a_i\in p\Z_p,\qquad
 h\in\mathcal O[[z]],\quad \bar h=\overline{h(0)}\ne0 .
\]
For \(u\ne0\), \(v(u)\ge b_r\), and primitive \(P\in\mathcal O[t]\)
with \(\deg\bar P\le q-2\), the cyclic coefficient vector of
\(\Pi_{n,q}(A)P(A)h\) modulo \(z^q-u\) is nonzero.
Here the component \(j\) of the vector of a series \(\sum f_kz^k\)
is \(\sum_{k\ge0}f_{qk+j}u^k\). These sums converge since \(v(u)>0\),
and their simultaneous vanishing is equivalent to vanishing at
the \(q\) distinct roots of \(z^q=u\).

At \(r=1\) reduce the vector modulo the maximal ideal.
The reduced ring is \(k[z]/(z^p)\), where \(k\) is the residue field,
\(h\) reduces to a constant unit, and \(\Pi_{n,p}(\bar A)=\bar A^n\).
The matrix of \(\bar A=z-zD+\bar c\) is triangular with distinct
diagonal entries \(\bar c-j\), \(0\le j<p\). They are all the
elements of \(\mathbb F_p\). Its vector \(1\) is cyclic because
\(\bar A^j1\) is monic of degree \(j\) for \(j<p\), so \(1\) has
a nonzero projection to each eigenspace. Vanishing would force
\(\bar P\) to vanish at all \(p-1\) nonzero eigenvalues, impossible
by its degree.

At \(r\ge2\) use the integral unit series
\[
 E(z)=\exp\!\left(\sum_{i\ge0}z^{p^i}/p^i\right)
     =\prod_{p\nmid m}(1-z^m)^{-\mu(m)/m},
\]
where \(\mu\) is the M\"obius function. The product identity follows
by taking formal logarithms and summing \(\mu\) over divisors.
Each factor and its inverse are integral because their exponents
are in \(\Z_p\) and the binomial coefficients are integral there
by continuity from nonnegative integers. The conjugation \(E^{-1}AE\)
subtracts \(zD\log E=\sum_{i\ge0}z^{p^i}\) from \(A\).
It preserves vanishing after changing the seed to \(E^{-1}h\),
since the series and their inverses converge at every point of the
fiber. In residue \(j\bmod p\), set \(y=z^p\),
\(b=a_1-1\), and \(w=by/p\). Write \(e_j\in\{0,\ldots,p-1\}\)
for the representative of \(c-j\bmod p\). The new block operator
is \(e_j+pL_j\), with
\[
 L_j=\frac{c-j-e_j}{p}-w\frac d{dw}+w+
        \sum_{i\ge1}\frac{(a_{i+1}-1)p^{p^i-1}}{b^{p^i}}w^{p^i}.
\]
This is in the same class. If \(E^{-1}h=\sum_{j=0}^{p-1}z^j
\widetilde h_j(y)\), the new seed
\(h_j(w)=\widetilde h_j(pw/b)\) reduces to the constant
\(\overline{h(0)}(-1)^j/j!\ne0\), using the first \(p\)
coefficients of \(e^{-z}\). Put \(q'=q/p\). The new fiber parameter
is \(u_*=(b/p)^{q'}u\), with \(v(u_*)\ge b_r-q'=b_{r-1}\).
If the original cyclic vector vanishes, the block vectors
\[
 \Pi_{n,q}(e_j+pL_j)P(e_j+pL_j)h_j
             \pmod{w^{q'}-u_*}
\]
all vanish. Indeed for each \(y^{{q'}}=u\), the evaluations on the
\(p\) roots of \(z^p=y\) separate these \(p\) residue blocks by a
Vandermonde matrix.

If \(e_j\ne0\), absorb the first factor into the seed \(\widehat h_j\).
This still reduces to a constant unit. The cyclic vectors of
\(L_j^i\widehat h_j\), \(0\le i<q'\), form a unit basis since
their reductions are polynomials of successively increasing degree
with unit leading coefficient in \(k[w]/(w^{q'})\).
The higher vectors are integral. Expanding \(P(e_j+pL_j)\)
in divided Taylor coefficients \(P^{[i]}(e_j)\) and inverting
the unit basis gives \(v(P^{[i]}(e_j))\ge q'-i>0\) for \(i<q'\):
all terms with \(i\ge q'\) are divisible by \(p^{q'}\).
Hence
\[
 \bar P(t)=(t^{p-1}-1)^{q'}H(t),\qquad \deg H\le q'-2.
\]
The first unit coefficient of \(P\) occurs at an index at most
\(q'-2\). Choose a coefficient \(\kappa\) of \(P(pt)\) of
least valuation, so \(v(\kappa)\le q'-2\).
Then \(Q(t)=P(pt)/\kappa\) is primitive and \(\deg\bar Q\le q'-2\).
In the unique block \(e_j=0\),
\(\Pi_{n,q}(pL_j)=p^n\Pi_{n,q'}(L_j)\), and division by
\(p^n\kappa\) contradicts the induction hypothesis at \(q'\).
This proves the invariant.

Finally the original \(\mathcal A,1\) belong to this class.
Scaling the dependent polynomial \(P\) to be primitive, the
rank condition contradicts the invariant. This proves
squarefreeness at every prime-power gap needed in this appendix.

\section{Intermediate terminal orders}\label{app:middle-tail}

\begin{proposition}\label{prop:middle-tail}
Assume the conditions of Lemma~\ref{lem:support-reduction}.
For \(2200\le N<e^{15}\), a common factor requires
\(\alpha<.03\) and \(b>N\). For \(3150\le N<e^{15}\)
there is no common factor.
\end{proposition}

Write \(t=\log(1/\alpha)\), \(V=\log(M+1)\) and
\(\sigma=2U+2\), where \(U=\log N\).
Throughout this appendix,
\(\log m=U+\log(1-1/N)>U-1/(N-1)>U-.001\).
From Lemma~\ref{lem:gap-caps} we may use \(d\le233\) below
\(20{,}000\) and \(d\le565\) throughout the appendix.
If \(X\) bounds the prime divisors of the lifted norm \(K\), the
right side of \eqref{eq:pade-norm} is at most
\begin{equation}\label{eq:intermediate-prime-head}
 S(X)+\pi(X)V/N,\qquad S(X)=\sum_{p\le X}\frac{\log p}{p-1}.
\end{equation}
In the applications below \(M\le X^j\), \(M\ge7N\) and
\(X\ge9\) give
\[
 \frac{V}{\log X}
 \le j+\frac{\log(1+1/M)}{\log X}
 <j+\frac{1}{7N\log9}<j+.001 .
\]
The baseline from Lemma~\ref{lem:resultant} is
\(\log\Lambda> U-1.011-g\). The weaker constant \(1.016\)
also suffices for the estimates of proper cosets.

\begin{lemma}\label{lem:small-prime-sums}
For each of the pairs
\[
 (x_0,c)=(21,.22),(34,.30),(43,.33),(100,.39),
          (117,.425),(137,.449),(150,.45),(343,.49)
\]
one has \(S(x)<\log x-c\) for \(x\ge x_0\).
For the prime count one has
\[
 \frac{\pi(x)\log x}{x}<
 \begin{cases}
  1.256 & (x\ge40),\\
  1.224 & (x\ge117),
 \end{cases}
\]
with the smaller bound on the overlap.
\end{lemma}

\begin{proof}
For \(S(x)-\log x\) below \(3000\), check the listed left endpoints
and the primes, since the function decreases between primes.
For the prime count, check \(40\) and \(117\) and the primes
up to \(300\), since between primes the multiplier \(\log x/x\)
decreases for \(x>e\). For \(x\ge300\),
Lemma~\ref{lem:prime-estimates} gives
\(1+1.2762/\log x\le1+1.2762/\log300<1.224\).
The supplement encloses the logarithms at all these checkpoints.
\end{proof}

\subsection{Proper rotation cosets}

Let \(k_0\) be the size of each rotation coset in
Lemma~\ref{lem:kummer}. Suppose \(q=d/k_0\ge2\), and put
\(P=M^{1/(q+1)}\). Then all norm primes are at most \(P\)
and \(4d\le2H<P\). In particular,
\begin{equation}\label{eq:proper-coset-two-caps}
 P^q<\min\{N/(4\alpha),N^2/4\}.
\end{equation}
The first inequality uses \(M=dN/\alpha<PN/(4\alpha)\),
and the second uses \(M<dN^2\).

If \(k_0=1\) and \(d\ge3\), then \(P^{d+1}=M<dN^2\).
Together with \(4d<P\) this gives \(4(4d)^d<N^2\).
In particular \(d\log(4d)<2U\), so \(d<U\) because
\(\log(4d)\ge\log12>2\).
Use \(X=(UN^2)^{1/4}>P\) and \(s<3U/4\) in the
baseline with \(g\le7s/32\). The resulting surplus over
\eqref{eq:intermediate-prime-head} is at least
\[
 U-1.016-\frac{21}{128}U
 -\log X+.22-4.001\cdot1.256\,X/N.
\]
Here \(\log X=U/2+(\log U)/4\).
The derivative of the terms linear or logarithmic in \(U\) is
\(1-21/128-1/2-1/(4U)>0\), and \(X/N\) decreases because
\(-1/2+1/(4U)<0\). The surplus therefore exceeds the certified
value \(1.10\) at \(N=2200\).
If \(k_0=1,d=2\), take \(X=(2N^2)^{1/3}\) and \(s<1.5\).
Now \(\log X=(2U+\log2)/3\), so the surplus again increases.
With the prime-sum correction \(.33\) in place of \(.22\),
it exceeds \(.95\) at \(N=2200\).

Suppose \(k_0\ge2\) and \(\alpha\ge.03\). For
\(\alpha\in[\ell,u]\) put \(X=\sqrt{N/(4\ell)}\).
By \eqref{eq:proper-coset-two-caps}, \(X>P\) and \(M<X^3\).
At least a third of each coset lies in the left half-plane.
The one-sided gain from Lemma~\ref{lem:resultant} gives the
lower bound
\[
 U-.006-\tfrac23(1.01+\kappa\sigma)
       +\tfrac13\log((1-u)/u)
\]
on \(\log\Lambda\), using \(g\le\kappa s\).
In every row \(X>40\), so \eqref{eq:intermediate-prime-head}
is bounded above by \(\log X+c+3.001\cdot1.256\,X/N\).
The lower bound on the gain has slope \(1-4\kappa/3>1/2\)
in \(U\). The logarithmic part \(\log X\) of the upper bound
has slope \(1/2\), and \(X/N\) decreases. It remains to check
the following margins at \(N=2200\).
\[
\begin{array}{cccc}
\toprule
{[\ell,u]}&\kappa&c&\text{surplus exceeds}\\
\midrule
{[.03,.1]}&.134&-.30&1.35\\
{[.1,1/3]}&.160&.20&.75\\
{[1/3,.75]}&.218&.20&.14\\
\bottomrule
\end{array}
\]

For \(\alpha<.03\) and \(q\ge3\), take
\(X=N^{2/3}/4^{1/3}>P\). From \(4d<P\) we obtain
\(M<dN^2<XN^2/4=X^4\).
If \(D\) bounds \(d\), the baseline exceeds
\(U-1.016-.00384D\).
Subtracting the upper bound
\(\log X-.22+4.001(X/N)(1+1.2762/\log X)\)
gives a surplus that increases for fixed \(D\), since
\(\log X\) has slope \(2/3\) in \(U\) and
\((X/N)(1+1.2762/\log X)\) decreases.
At \(N=2200,D=233\) and \(N=20{,}000,D=565\) it
exceeds \(1.08\) and \(.68\), respectively.

For \(\alpha<.03\) and \(q=2\), use \(X=N/2\), \(M<X^3\)
and \eqref{eq:intermediate-prime-head} to bound the norm
charge by \(U+.862\).
When \(d=4\), half of the size-two coset has gain at least
\(m(1-\alpha)/\alpha\), and the mean log gain exceeds
\(U+1.21\). For \(6\le d\le18\), the small-coset form of
Lemma~\ref{lem:angles}, with \(k_0=d/2\), gives a gain greater
than \(U+1.89\). For \(d\ge20\), use \(g_0\le3.18\),
\(g_0/k_0\le.10868\), \(r/k_0<.111\) and \(J<.153\)
in \eqref{eq:pade-angles}, together with
\(\log C'>t-.7237\) and \(1-1/k_0\ge.9\).
Dropping the positive \(2\log(2k_0/\pi)/k_0\) leaves a gain
exceeding \(U+1.80\).
This proves that the rotation coset is full.

\subsection{Full cosets at moderate and large ratios}

First take \(2200\le N<20{,}000\).
Here \(a\le(97/3)d\), \(d\le233\) and \(H\le467\).
Thus \(aN+H<189{,}695{,}659=X_{222}\).
Lemma~\ref{lem:local-gap} gives \(H<222\) and \(d\le110\).
For \(N\le3500\) these caps give \(aN+H<X_{154}=17{,}051{,}707\),
so \(d\le76\).

Let \(\ell\le\alpha\le u\) and order the \(d\) rotations by
increasing absolute argument, starting at index zero. An arc
centered at angle zero and containing \(j+1\) arguments has
length at least \(2\pi j/d\), by the rotation spacing.
Thus \(c_j=\cos(\pi j/d)\) bounds the cosine in position \(j\).

For a nonnegative upper bound \(c\) on the cosine, set
\(W=(1-u-c)/u\). In the second integrand in
\eqref{eq:pade-integrals}, write \(v=1-e^{-x/d}\).
Use \(v\le x/d\) and \(\Re z/m\le bc\) to bound the logarithm
of the original integrand, divided by \(m\), by
\[
 -\frac{1-\alpha-c}{\alpha}x+\log(1-e^{-x})
 \le -Wx+\log(1-e^{-x}).
\]
For \(W>0\) the derivative on the right vanishes only at
\(e^{-x}=W/(1+W)\), and the second derivative is negative.
The maximum is \(-J(W)\), where
\(J(W)=(1+W)\log(1+W)-W\log W\).
The original integral over \(0\le v\le1\) is at most
\(e^{-mJ(W)}\). Since \(J(W)<1+\log(1+W)\le
1+\log(100/3)<5\), Stirling's inequality gives
\((m!e^{mJ(W)})^{1/(m+1)}>me^{-1.01+J(W)}\).
Indeed \(.01(m+1)+\tfrac12\log(2\pi/m)-5\) is positive at
\(m=2199\) and has derivative \(.01-1/(2m)>0\).
The linear gain in Lemma~\ref{lem:resultant} is at least \(mW\).
Together the two gains give
\(\Phi_u(c)=\max\{-1.01+J(W),\log W\}\) after subtracting
\(\log m\).
For \(c<0\) use \(W=(1-u)/u\) in the integrand and
replace \(\log W\) in this maximum by \(\log(W(1-c))\).
For \(c\ge0,W\le0\) use the baseline
\(-1.01-\kappa(u)du\), where \(\kappa(u)=g/s\) at \(\alpha=u\)
in \eqref{eq:pade-loss}.
In the smaller terminal band we also use the lower bound
\[
 -1.01-\frac d u
       \left(c-(1-u)-(1-u)\log\frac{c}{1-u}\right)
\]
when this improves the baseline. It follows by discarding
the beta tail factor and maximizing the remaining exponent.
The loss at \(\alpha\) is zero for \(c\le1-\alpha\) and
otherwise the same formula with \(\alpha\) in place of \(u\).
For \(c>1-\alpha\) the derivative of the loss is
\[
 \frac d{\alpha^2}
 \left(\log\frac{c}{1-\alpha}-c+1-\alpha\right)\ge0.
\]
For the inequality, \(\log(c/(1-\alpha))\) is the integral
of \(1/x\ge1\) from \(1-\alpha\) to \(c\le1\).
This permits evaluating the loss at the upper endpoint \(u\).

By Lemma~\ref{lem:degree-cap}, a prime bound is
\(\max(a+2d-1,\sqrt M)\). It lies between \(X_1\) and \(X_0\),
where
\[
 X_0=\max\{d(1+1/\ell),\sqrt{dN/\ell}\},\qquad
 X_1=\sqrt{\max(7,d/u)N}.
\]
The upper bound for the right side of the norm inequality is therefore
\begin{equation}\label{eq:moderate-norm}
 \log X_0-c_0+2.001(1+1.2762/\log X_1)X_0/N.
\end{equation}
Use \(c_0=.39\) in the first band, where \(X_1>100\), and
\(c_0=.45\) in the second, where \(X_1>150\).
In \(U\), \(\log X_0\) is the maximum of a constant and an
affine function of slope \(1/2\). The quotient \(X_0/N\) is
the maximum of two decreasing functions, and \(1/\log X_1\)
also decreases. Thus the norm bound grows at most at rate
\(1/2\). The angular lower bound
\(\log(N-1)+d^{-1}\sum_j\Phi_u(c_j)\) grows at rate \(N/(N-1)>1\).
Since \(d\ell\le s<2U+2\), the first eligible order is
bounded below by \(\max\{N_{\min},\exp((d\ell-2)/2)\}\).
For each \(d\) and each interval of the following subdivision,
the interval checker evaluates the surplus at that lower bound.
Intervals with \(d(1/\ell-1)<5\) contain no eligible support.
\[
\begin{array}{rrlc}
\toprule
N_{\min}&d_{\max}&\text{interval widths}&\text{surplus exceeds}\\
\midrule
2200 &76 & .005\text{ to } .05,\ \text{then } .01 &.070\\
3500 &110& .02 &.089\\
\bottomrule
\end{array}
\]

Now take \(20{,}000\le N<e^{15}\). The gap table gives the
successive caps \(D=257,300,565\) at starting orders
\(20{,}000,32{,}768,65{,}536\). In each range \(D/N<.015\).
Since \((1+\alpha)^2/\alpha\) decreases on \([.03,.75]\),
\[
 \frac{(1+\alpha)^2b}{N}
 \le\frac D N\frac{(1+\alpha)^2}{\alpha}
 <.015\frac{1.03^2}{.03}<1.
\]
Hence \(\sqrt M\) bounds the norm primes. Using
\(S(\sqrt M)<\log\sqrt M-.49\), the right side of the norm inequality is below
\begin{equation}\label{eq:moderate-high-norm}
 \tfrac12(U+\log b)-.49+2.46\sqrt{b/N}.
\end{equation}
At angles with \(\cos\theta\le1-\alpha\) the beta integral
has gain at least \(me^{-1.01}\). The exceptional fraction
is at most \(\arccos(1-\alpha)/\pi+1/d\). For
\(\ell\le\alpha\le u\), use \(\eta\ge\arccos(1-u)/\pi\).
At \(d\ge10\), at least a fraction \(.45\) lies in the left
half-plane, where one may use the greater gain \(m(1-\alpha)/\alpha\).
Thus, with \(S_0=\min(2U+2,Du)\) and
\(b_0=\min((2U+2)/\ell^2,D/\ell)\), the surplus exceeds
\[
 U-1.011-\kappa(\eta S_0+u)
  +.45\max\{0,1.01+\log((1-u)/u)\}
  -\bigl(\tfrac12(U+\log b_0)-.49+2.46\sqrt{b_0/N}\bigr).
\]
Where \(S_0\) is differentiable in \(U\), its derivative is
zero or two. Thus the lower bound on the gain grows at rate
at least \(1-2\kappa\eta\). The logarithmic term in the
upper bound grows at most at rate \(1/2+1/(2U+2)\), since
\((\log b_0)'\le2/(2U+2)<1\). The latter inequality also
makes \(\sqrt{b_0/N}\) decrease. The checker verifies
\(1-2\kappa\eta>1/2+1/(2U+2)\) at each starting order,
where the comparison is tightest. It also checks
\(\cos(\pi\eta)<1-u\) and the upper bound \(\kappa\)
on \(g/s\) at \(u\). The table records the least certified
surplus across the three starting orders for each interval.
For \(d\le9\) the same argument before adding the left half-plane
gain uses \(\kappa=.218,\eta=.42,\ s\le6.75,\ b\le300\),
and gives a positive surplus at \(N=20{,}000\) increasing thereafter.
\[
\begin{array}{cccc}
\toprule
{[\ell,u]}&\kappa&\eta&\text{surplus exceeds}\\
\midrule
{[.03,.036]}&.128&.086&.12\\
{[.036,.045]}&.129&.096&.20\\
{[.045,.07]}&.131&.120&.11\\
{[.07,.11]}&.134&.151&.21\\
{[.11,.17]}&.140&.189&.51\\
{[.17,.28]}&.150&.245&.67\\
{[.28,.45]}&.170&.315&.62\\
{[.45,.60]}&.190&.370&.53\\
{[.60,.75]}&.218&.420&.08\\
\bottomrule
\end{array}
\]


\subsection{Small ratios with the upper exponent at most the order}

Take \(X=1.03\sqrt M\), which is a prime bound by
Lemma~\ref{lem:degree-cap}. Indeed
\(a+2d-1<(1+\alpha)b\le1.03\sqrt M\).
For \(d\ge20\), the angular estimate
\eqref{eq:pade-angles} gives
\[
 \log\Lambda-U>
 .95\bigl(\log(100/3)-.7237\bigr)-.002-.05434
 - .111\bigl(1.01+.00384D+\log2\bigr)-.153
\]
with \(d\le D\). In \eqref{eq:intermediate-prime-head},
use \(X/N\le1.03\) and \(X>1.03\sqrt{20N/.03}\).
At \(N=2200,D=233\) the surplus exceeds \(.17\), and at
\(N=20{,}000,D=565\) it exceeds \(.08\).
For fixed \(D\), the gain bound minus \(U\) is constant.
After using \(X/N\le1.03\) in the norm bound, its remaining
\(N\)-dependence is through the decreasing reciprocal of
\(\log(1.03\sqrt{20N/.03})\).
For \(10\le d\le19\), inserting each \(d\) in
\eqref{eq:pade-angles} gives \(\log\Lambda>U+2.41\),
against a right side of the norm inequality below \(U+2.052\) by \(X>343\).
For \(4\le d\le9\) use the small-coset angular bound. The gain
exceeds \(U+2.24\), against a right side of the norm inequality below \(U+2.177\)
by \(X>137\).

For an antipodal coset with \(d=2\), write \(c=|\cos\theta|\),
\(A_0=(1-\alpha)/\alpha\) and \(B_0=e^{-g_0}\).
The two gains have product at least
\(m^2A_0(1+c)\max\{B_0,A_0-c/\alpha\}\).
On \(0\le c\le1-\alpha(1+B_0)\) this is a concave quadratic
in \(c\) and is linear increasing on the rest of its domain.
The two endpoint values in the quadratic range, divided by
\(m^2A_0\), are \(A_0\) and \(B_0(2-\alpha(1+B_0))\).
The second is smaller because \(A_0>32\) and \(B_0<1\).
This gives the lower bound \(m^2A_0B_0(2-\alpha(1+B_0))\).
In particular \(\log\Lambda>U+t/2-.2\).
For \(d=3\) the small-coset bound gives
\(U+2t/3-.435>U+t/2-.2\).
Put \(\mu=\tfrac12\log(b/N)\).
Since \(X>117\), the right side of the norm inequality is at most
\[
 U+\mu+\log(1.03)-.425+
       2.001\cdot1.03\cdot1.224\,e^\mu.
\]
Substitute \(t=U-\log d+2\mu\) in the gain \(U+t/2-.2\).
The linear terms in \(\mu\) cancel in the surplus, while the
remaining negative multiple of \(e^\mu\) is smallest at
\(\mu=0\). The resulting bound increases with \(U\)
and exceeds \(.96\) at \(N=2200\).
This proves the first conclusion of Proposition~\ref{prop:middle-tail}.

\subsection{Small ratios with the upper exponent exceeding the order}

It remains to exclude this case for \(N\ge3150\).
The bounds \(H<G,\ d\le D=G/2-1\) are shown in the
following table. In addition to the main gap table, apply
Lemma~\ref{lem:local-gap} successively using \(G=320,292,288,282\)
up to \(N=5000,4150,3467,3295\), respectively.
At every step the preceding caps give
\(aN+H<(5/6)D_{\mathrm{old}}N_{\max}^2+H_{\mathrm{old}}<X_G\).
\[
\begin{array}{rrrccc}
\toprule
N\text{ from}&N\text{ through}&G&A_v&A_c&\text{surplus exceeds}\\
\midrule
3150  &3295   &282 &2.504&2.885&.023\\
3296  &3467   &288 &2.506&2.886&.146\\
3468  &4150   &292 &2.514&2.889&.31\\
4151  &5000   &320 &2.522&2.892&.44\\
5001  &6000   &354 &2.531&2.896&.86\\
6001  &10000  &456 &2.58 &3.01 &.59\\
10001 &19999  &468 &2.58 &2.92 &2.37\\
20000 &32767  &516 &3&3&3.88\\
32768 &65535  &602 &3&3&4.66\\
65536 &3269017&1132&3&3&4.87\\
\bottomrule
\end{array}
\]

Let
\[
 \begin{aligned}
 w&=\log(b/N), & v&=\log M=2U+w, & P&=4M^{1/3},\\
 \zeta&=\log P, & t_D&=U+w-\log D, & \xi&=t_D+.337.
 \end{aligned}
\]
We have \(0<w<\log(5D/6)\) and \(P>\max(830,5G/2)\)
in every row. At \(p>P\) we have \(p>2H+d\) and \(p^3>M\).
Lemma~\ref{lem:two-heights} separates the prime divisors of the
norm there according to whether the two Newton edge heights agree.

The primes with equal heights divide
\(\binom{a+2d-1}{2d-1}\), since they divide a number in
\([a+1,a+2d-1]\) and exceed \(2d-1\).
Using \(\binom{x}{j}\le(ex/j)^j\), their total logarithmic weight
is at most
\[
 2d\bigl(1+\log(a/d+2)-\log2\bigr)<2d(t+.337)\le2D\xi.
\]
For the last step the function \(x(\log(b/x)+.337)\) increases
for \(2\le x\le D\), since \(b>N\) and \(D<N/4\).
Their right side of the norm inequality is at most
\(2D\xi(1/(P-1)+V/(N\zeta))\).

For unequal heights, one of the two windows
\([bm-2H,bm]\), \(m=N-1,N\), contains \(qp^2\)
with \(1\le q\le Q=M/P^2=P/64\).
At fixed \(q\) each window contains at most one such \(p>P\),
since consecutive integer squares multiplied by \(q\) are
spaced by more than \(2P>2H\) there.
Thus the digit charge is at most \(2QV/N\).
Put \(Z=b(N-1)-2H\). Since \(Z>.9995M\) and \(Q>12\),
the first sum in \eqref{eq:pade-norm} on these primes is at most
\[
 v\,\frac{\sum_{q\le Q}\sqrt q}{\sqrt Z-\sqrt Q}
 \le \frac{1.002v}{64}\left(\frac23+\frac{1}{2Q}\right)
 < .71v/64.
\]
Here \(p\ge\sqrt{Z/q}\) and \(p\le\sqrt M\).
For the middle bound use the trapezoid inequality for the concave
square root and
\(\sqrt Z-\sqrt Q\ge\sqrt M(.9995-1/P)>\sqrt M/1.002\).
The lower bound on \(Z/M\) follows from
\(1-1/N-2H/N^2>.9995\), with \(H\le467\) below
\(20{,}000\) and \(H\le1131\) above.

The primes \(p\le P\) contribute at most
\(S(P)+\pi(P)V/N\). We have \(V<v+.001\) and
\[
 \frac{V}{\zeta}<3-\frac{3\log4-.001}{\zeta}.
\]
Put \(k=3\log4-.001\) and \(c_1=1.2762\).
Both \(3-k/\zeta\) and
\[
 (3-k/\zeta)(1+c_1/\zeta)
 =3-\frac{k-3c_1}{\zeta}-\frac{kc_1}{\zeta^2}
\]
increase in \(\zeta\), since \(k>3c_1\).
This also bounds both by three. In each row the upper corner
has its largest \(N\) and \(b/N=5D/6\), so the checker
evaluates there to obtain \(A_v,A_c\).
Combining the three ranges of primes, using \(S(P)<\log P-.49\)
and \(P/(P-1)<1.002\), gives the upper bound
\begin{equation}\label{eq:hard-upper}
 R=\zeta-.49+\frac PN\left(A_c+\frac{2(v+.001)}{64}\right)
 +\frac{.71v}{64}+\frac{2.004D\xi}{P}+\frac{2A_vD\xi}{N}.
\end{equation}

For the lower bound let \(E=D/b\).
The errors in \eqref{eq:pade-angles} satisfy
\(J<.883\sqrt\alpha\) and \(r/d<.641\sqrt\alpha\).
Indeed the derivative of \(\arcsin x+x\sqrt{1-x^2}\)
is \(2\sqrt{1-x^2}\le2\). For the second estimate,
convexity bounds \((1-y)^{-1/2}\) on \(0\le y\le.03\)
by the chord from \(1\) to \(1/\sqrt{.97}<1.0156\).
Integrating the resulting upper bound \(1+.52x^2\) gives
\(\arcsin x\le x+.52x^3/3<1.006x\) for \(0<x^2\le.03\).
For \(d\ge20\), use
\(\log C'>t-.6932-1.031E\),
\(g_0\le1.01+.128DE\), and \(g_0/d\le1.01/d+.128E\).
At fixed \(b\), the derivatives of
\((1-1/x)(\log(b/x)-.6932)\) and
\([2\log(2x/\pi)-1.01]/x\), multiplied by \(x^2\), are
\[
 \log(b/x)-.6932-x+1,\qquad 3.01-2\log(2x/\pi),
\]
respectively. Both are negative for \(d\le x\le D\), since
\(x\ge20\) and \(\log(b/x)\le t<U<15\).
Substituting \(D\) for \(d\) in these two decreasing terms
and bounding the errors with \(E\) gives
\begin{equation}\label{eq:hard-lower}
 \begin{aligned}
 L={}&U-.001+(1-1/D)(t_D-.6932)
     +\frac{2\log(2D/\pi)-1.01}{D}\\
     &-1.159E-\bigl[.641(1.7032+.128DE)+.883\bigr]\sqrt E .
 \end{aligned}
\end{equation}
For \(d=2,3,\ 4\le d\le9,\ 10\le d\le19\), use
\(D=2,3,9,19\) respectively in \eqref{eq:hard-upper} and
\(t_D\), keeping the row's \(A_v,A_c\).
The same angular estimates used above give
\(L=U+ft_D-c\) with
\((f,c)=(1/2,.2),(2/3,.435),(.75,.68),(.9,.925)\),
respectively. For \(10\le d\le19\) retain the positive term
\(2\log(2d/\pi)/d\ge2\log(38/\pi)/19\) in
\eqref{eq:pade-angles}.

For fixed \(U\), the lower bound \(L\) is affine in \(w\)
minus a nonnegative linear combination of \(e^{-w}\),
\(e^{-w/2}\) and \(e^{-3w/2}\). It is therefore concave.
For \(R\) only the \(P/N\) and \(\xi/P\) terms have nonzero
second derivatives in \(w\). Their sum gives
\[
 R_{ww}=
 \frac{P}{9N}\left(A_c+\frac{2(v+.001+6)}{64}\right)
      +\frac{2.004D}{9P}(\xi-6).
\]
To prove \(R_{ww}\ge0\), multiply by \(9P\) and compare the
positive term with \(2.004D(6-\xi)\).
The positive term increases in \(U,w\) because
\(P^2/N=16e^{U/3+2w/3}\) and \(v=2U+w\).
The term \(D(6-\xi)\) decreases in \(U,w\).
Where it is positive, it increases with \(D\) because its
derivative in \(D\) is \(7-\xi>1\).
It suffices to check \(w=0\) at the smallest order in each row
with its largest \(D\).

For fixed \(w\), \(L\) grows with \(U\) at rate at least \(1.5\),
since its negative exponential terms increase towards zero.
The \(U\)-derivative of \((P/N)(A_c+(v+.001)/32)\) equals
\[
 \frac PN\left(\frac1{16}-\frac{A_c+(v+.001)/32}{3}\right)<0,
\]
since \(A_c>2\).
The two equal-height terms decrease because their derivatives
have the signs of \(1-2\xi/3\) and \(1-\xi\), and \(\xi>3\).
Thus \(R\) grows at rate at most \(2/3+2\cdot .71/64\).
The supplement checks the convexity inequality and the surplus
at the two endpoints of \(w\) at the smallest \(N\) in each row.
For \(d\ge20\) the certified margins are given in the table.
For all four smaller-gap cases the margin exceeds \(3.8\).
It follows that \(L>R\) throughout each row, which contradicts
\eqref{eq:pade-norm} and completes the proposition.


\section{The range \texorpdfstring{\(N\ge e^{15}\)}{N at least exp(15)}}\label{sec:large-tail}

\begin{proposition}\label{prop:large-tail}
Suppose \(U=\log N\ge15\), \(a,d\ge1\), \(d<3a\), and
\[
  s=\frac{d^2}{a+d}<2U+2,\qquad 6a<(5N-11)d .
\]
Then \(R_{N-1}\) and \(R_N\) are coprime.
\end{proposition}

\begin{proof}
Suppose there is a monic irreducible common factor \(F\).
The notation \(m,b,M,h,H,\alpha,V,\Lambda\) is as in
Lemma~\ref{lem:resultant}. Set \(\sigma=2U+2\), \(t=\log(1/\alpha)\),
and \(B=H^4/M\). The starting-point inequality of
Lemma~\ref{lem:degree-cap} is available throughout this range,
and \(\log m>U-.001\).
The shape inequality gives \(b<(5/6)dN\), hence \(t<U\)
and \(d<\sigma N\). Thus \(M<\sigma e^{3U}\).
The comparison \(e^{.24U}>\sigma+e^{-3U}\) at \(U=15\)
extends to larger \(U\) because \(.24e^{.24U}>2\) there.
It follows that \(V<3.24U\).
We first force the full set of rotations in the Padé norm.
The angular gain then rules out small \(H\).
For the remaining factors, counts of square multiples in short
intervals bound the large-prime contribution to the norm.

\subsection*{Rotations and the range of the parameters}

We record how the weak bound in Lemma~\ref{lem:angles} will be used.
For a coset of size \(n\) and \(\alpha<.03\), put
\(c_\alpha=\log(4/(1-\alpha))\). Discarding the positive
\(2\log(2n/\pi)/n\) term and the nonnegative term
\(-(r/n)\log(1-\alpha)\) gives the lower bound
\begin{equation}\label{eq:large-angular-template}
 U-.001+(1-1/n)(t-c_\alpha)
               -g_0/n-(r/n)g_0 .
\end{equation}
Here \(c_\alpha<1.417<t\) and \(g_0<1.266+.256U\).
For \(d\ge2\) and \(s<.27\), a full coset also gives
\(\log\Lambda>U+t/3-.71\): at least a third of the rotations
have log gain at least \(\log m+t+\log(1-\alpha)\), and the
baseline gain at every rotation exceeds \(\log m-1.01-.128s\).
Use \(s<.27\) to bound the resulting loss.

Suppose first that \(q=d/k_0\ge2\) in Lemma~\ref{lem:kummer},
and set \(P=M^{1/(q+1)}\). The right side of
\eqref{eq:pade-norm} is bounded by \(S(X)+\pi(X)V/N\) for any
\(X\ge P\). If \(\alpha\in[\ell,u]\subseteq[.03,.75]\),
the inequality \(M=dN/\alpha<PN/(4\alpha)\) gives
\(P^q<N/(4\alpha)\). Take \(X=(N/(4\ell))^{1/2}\), so
\(X>P\), \(M<X^3\), and \(V/\log X<3.001\).
Here \(X\ge\sqrt{3N/4}\) and
\(\log(M+1)\le\log M+1/M\). Thus
\(\pi(X)V/N<3.6X/N\), since
\(3.001(1+1.2762/\log\sqrt{3N/4})<3.6\) at \(U=15\).
Using \(S(X)<\log X-.49\), the surplus of the baseline
\(\log m-1.01-g\) over the upper bound is greater than
\[
 \tfrac12(U+\log(4\ell))-1.011-\kappa(2U+2)+.49
                  -\frac{3.6}{\sqrt{4\ell N}},
\]
where \(g\le\kappa s\). Use \(\kappa=.134,.160,7/32\),
respectively, for \([.03,.1]\), \([.1,1/3]\), and
\([1/3,.75]\), by \eqref{eq:pade-loss}.
The smallest endpoint surplus at \(U=15\) is greater than
\(.12\). The slopes of the affine terms are
\(\tfrac12-2\kappa\ge1/16>0\), and the last (subtracted)
term decreases.

For \(\alpha<.03\) and \(q\ge3\), the inequalities
\(M<dN^2<PN^2/4\) give \(P<X=N^{2/3}/4^{1/3}\)
and \(M<X^4\). The baseline exceeds \(.744U-1.267\).
Here \(V/\log X<4.001\) and
\(4.001(1+1.2762/\log X)<4.6\) at \(U=15\), so the norm
bound is \(\log X-.54+4.6X/N\).
The surplus exceeds \(.85\) at \(U=15\). The baseline grows
at rate \(.744>2/3\) and \(X/N\) decreases.
If \(\alpha<.03\) and \(q=2\), we have \(P/N<1/2\) and
\(V/\log P<3.001\). The digit loss \(\pi(P)V/N\) is less
than \(4.848P/N<2.43\), since \(P>2H\ge8\).
For \(d<50\) this improves to \(.121\) using
\(P/N<(50/N)^{1/3}\). In that case the gain exceeds \(U-1.22\)
and the prime sum is below
\(\max\{8,(2U+\log50)/3\}-.54\). For \(d\ge50\), the
weaker bound in Lemma~\ref{lem:angles}, with \(n=d/2\),
gives gain greater than \(.9715U+.96t-1.608\) in logarithm.
In \eqref{eq:large-angular-template}, use \(1/n\le.04\),
\(g_0/n\le1.01/25+2(.128)(.03)<.049\), and \(r/n<.111\).
The prime sum and digit loss are less than
\(\max\{8,(U+2t+\log\sigma)/3\}+1.89\).
The surplus is positive at \(U=15,t=3.5\).
Its two branches increase with \(U\) and \(t\), since the
derivatives of \((U+2t+\log\sigma)/3\) in these variables
are at most \(17/48\) and \(2/3\), respectively. We conclude that \(k_0=k=d\).

For this full coset, suppose \(\alpha\ge.03\).
If \(b=2\), then \(a=d=1\), and the norm primes are at
most \(\sqrt{2N}\). The prime-count bound and
\(S(\sqrt{2N})<\log\sqrt{2N}-.49\) already contradict
the baseline \(U-1.011-.176/2\).
Otherwise \(b\ge3\), so \(\log\sqrt M>8\). Since
\(b\le\sigma/(.03)^2<N/4\), the degree and prime bounds
give prime support at most \(\sqrt M\). Here \(V/\log\sqrt M<2.001\) and
\(2.001(1+1.2762/(U/2))<2.35\) at \(U=15\).
The right side of \eqref{eq:pade-norm} is then less than
\begin{equation}\label{eq:large-moderate-upper}
 \tfrac12(U+\log b)-.54+2.35\sqrt{b/N}.
\end{equation}
For \(d\ge10\), at least \(.45\) of the rotations lie in the
left half-plane. On the first and last intervals
\([\ell,u]\) in the table below, the log gain therefore exceeds
\(U-.001-.55(1.01+\kappa\sigma)+.45\log((1-u)/u)\).
On the middle interval at least \(17/30\) of the rotations
have \(|\arg z|\ge\pi/3\), hence \(\Re z\le am\). There the
gain is at least \(me^{-1.01}\), so the logarithmic mean
exceeds \(U-1.011-(13/30)\kappa\sigma\).
Using \(b\le\sigma/\ell^2\), the limiting comparisons are
\[
\begin{array}{c|c|cc}
{}[\ell,u]&\kappa&\text{gain lower bound at }U=15
                    &\text{upper bound at }U=15\\ \hline
{}[.03,.1]&.134&13.07&12.45\\
{}[.1,.5]&.176&11.54&11.08\\
{}[.5,.75]&7/32&10.09&9.41
\end{array}
\]
The affine lower bounds have slope at least \(.75\), whereas
the logarithmic part of the upper bounds has slope at most
\(1/2+1/\sigma\) and the square-root term decreases.
For \(d\le9\),
use \(s<6.75,\ b\le300\) and the baseline
\(U-1.011-7s/32\) against \eqref{eq:large-moderate-upper}.
Hence \(\alpha<.03\) is necessary.

If \(b\le N\), the norm primes are at most \(X=1.03\sqrt M\).
Use \(\log X>8\), \(V/\log X<2.001\), and
\(\pi(X)V/N<2.42\sqrt{b/N}\).
The right side of \eqref{eq:pade-norm} is less than
\(\tfrac12(U+\log b)+1.91\), since
\(\log(1.03)-.54+2.42<1.91\).
When \(d\ge10\), use \(n=d\) in
\eqref{eq:large-angular-template} with \(1/n+r/n<.211\).
This gives \(\log\Lambda>.945U+.9t-1.70\). Writing
\(.9t-\tfrac12\log b=-.05\log b-.45\log s\) and using
\(\log b\le U\), the surplus is greater than
\(.395U-.45\log\sigma-3.61>0\).
This last expression increases, since \(2/\sigma\le1/16\).
If \(2\le d\le9\), then \(s=d\alpha<.27\), so the gain
\(U+t/3-.71\) gives a surplus
greater than \(U/3-(\log9)/2-2.62>0\), using
\(\log b=t+\log d\) and \(t<U\).
For \(d=1\), the last gain in Lemma~\ref{lem:resultant}
exceeds \(U+\tfrac12\log a-1.195\) in logarithm.
Against \(\tfrac12(U+\log b)+1.91\) this leaves a positive
margin, using \(b=a+1\le2a\). The shape inequality precludes
\(b>N\) when \(d=1\), so now \(b>N\) and \(d\ge2\).

For any \(x\ge\max(e^8,\log K)\), the prime sum in
\eqref{eq:pade-norm} is below \(\log x+.461\). Below \(x\)
use Lemma~\ref{lem:prime-estimates}, and above \(x\) bound it
by \(\log K/(x-1)\le x/(x-1)<1.001\).
Suppose \(B<2048\). Use \(B=h^4s^2b/N\) for the bound on \(s\). For the bound
on \(V\), write \(M=(B/h^4)^{1/3}e^{(4t+4U)/3}\) and
use \(h\ge2\) and \(t<U\). These give
\[
 s<8\sqrt2,\qquad V<\tfrac83U+1.62,\qquad
 H<10.08e^{(U+t)/3}.
\]
The last bound follows from \(H=(B/h)^{1/3}e^{(U+t)/3}\).
Split the count of distinct norm primes at \(N/16\)
and use \(\log K\le HV\). The digit loss is less than
\(.236+9.46UH/N\). Both coefficients follow from the prime count at \(U=15\).
To see the decrease, divide the bound \(\tfrac83U+1.62\)
for \(V\) by \(U-\log16\) in the first coefficient, and use
\((\tfrac83+1.62/U)^2/(1-\log16/U)\) in the second.
For \(d\ge50\) the bound \eqref{eq:large-angular-template} gives
\(\log\Lambda>U+.98t-1.72\), using \(g_0<2.459\),
\(g_0/d<.025\), and \(r/d<.111\).
Splitting the prime sum at \(\max(e^8,HV)\), the surplus is
greater than
\[
 \tfrac23 U+(.98-\tfrac13)t
 -\log(10.08(\tfrac83 U+1.62))-2.417
 -95.3568\,Ue^{(t-2U)/3}.
\]
This is concave in \(t\). At \(t=3.5,U\) its values are
greater than \(3.59,1.60\), respectively, at \(U=15\).
At either endpoint the logarithmic term grows at rate less
than \(1/U\), while \(Ue^{-jU/3}\) decreases for \(j=1,2\).
The endpoint surpluses therefore increase. If \(d<50\), then \(s<2500/N<.001\)
and \(t>U-\log50\). The gain exceeds \(U+t/3-.71\).
Use \(M<50N^2\), \(H<(102400)^{1/4}N^{1/2}\), and
\(V<2U+\log51\). The digit bound is below \(1.65\), so
the total upper bound is less than
\[
 \tfrac12U+\tfrac14\log102400+\log(2U+\log51)+.461+1.65.
\]
It is beaten at \(U=15\) by the gain, which increases faster.
Thus \(B\ge2048\).

Finally, the initial block at order \(N\) has prime support
below \(1.03b\), since \(b>N\) and \(\alpha<.03\).
Its logarithm exceeds \(H\log(aN)\). For each \(p^j\le M\), the initial block contains at most
\(H/p^j+1\) multiples of \(p^j\). Summing over \(j\) and
then over the supported primes bounds its logarithm by
\(H S(1.03b)+\pi(1.03b)\log M\).
The difference \(\log(aN)-S(1.03b)\) exceeds \(U+.47\)
and the prime-count term is below \(2.24b\). We have obtained
\begin{equation}\label{eq:large-parameter-range}
 k_0=k=d,\qquad \alpha<.03,\qquad b>N,\qquad
 B\ge2048,\qquad H/b<\frac{2.24}{U+.47}.
\end{equation}

\subsection*{Norm and degree bounds}

Put \(P=2H+d\) and \(L=\log M\). For primes \(p>P\) one
has \(p^4>M\). Call a prime exceptional if either terminal
interval of length \(6H\) contains a positive multiple of
\(p^3\). For each such interval and each \(q\le M/P^3\)
there is at most one prime with \(qp^3\) in the interval.
Each exceptional prime contributes at most \(L\) to the log
norm by \eqref{eq:norm-affine}. Write \(Y\) for the sum of
\(\log p\) over the nonexceptional unequal-height norm primes.
By Lemma~\ref{lem:two-heights} and \eqref{eq:unequal-products},
\begin{equation}\label{eq:large-two-y}
 Y\le \min\{6d\log(2M),\,ML/P^2\}.
\end{equation}
The second bound uses at most one prime per square multiplier
in each terminal interval and \(\log p\le L/2\).

At a nonexceptional equal-height prime its entire norm
power divides an integer in \([a+1,a+2d-1]\).
The product of these powers divides \(\binom{a+2d-1}{2d-1}\)
because the primes exceed \(2d-1\).
Writing \(A_0=a/d\), its logarithm is less than
\(2d(1+\log(A_0+2)-\log2)\).
At \(p\le P\) apply the affine estimate
\eqref{eq:norm-affine}. Comparing with
\(\log|F(0)|>H\log(aN)\), and setting
\[
 \Delta=\log(aN/H)-\frac2h(1+\log(A_0+2)),
 \qquad \Delta_*=\Delta+\frac2h\log2 ,
\]
gives
\begin{equation}\label{eq:large-norm}
 \Delta_*<\log(P/H)+\frac{\pi(P)L}{H}+\frac{Y}{H}+\frac{L}{4B}.
\end{equation}
Here \(\log P>8\), so \(S(P)<\log P\), and
\(2ML/(HP^3)\le L/(4B)\). The same estimate holds for \(\Delta\).

In \eqref{eq:large-parameter-range} we have
\(A_0+1>e^{U/2}/\sqrt{2U+2}>308\).
For \(h=2\), \(\Delta-U=\log(A_0/(A_0+2))-1-\log2>-1.70\),
and \(\Delta_*>U-1.01\).
For \(h\ge4\), \(\Delta-U\) attains its minimum on
\(4\le h\le 2.24(A_0+1)/(U+.47)\) at an endpoint, since
its derivative has the sign of \(2(1+\log(A_0+2))-h\).
At the lower endpoint the value exceeds \(.97\).
At the upper endpoint it is
\[
 \log\frac{U+.47}{2.24}+\log\frac{A_0}{A_0+1}
 -\frac{2(U+.47)}{2.24}\frac{1+\log(A_0+2)}{A_0+1}.
\]
This increases with \(A_0\). Substitute the lower bound
\(A_0+1>e^{U/2}/\sqrt{2U+2}\). The result exceeds \(1.6\)
at \(U=15\) and increases thereafter. Indeed, if
\(T=e^{U/2}/\sqrt{2U+2}\), then \(.46<T'/T<.5\).
The logarithmic derivative of
\((U+.47)(1+\log(T+1))/T\) is at most
\(1/(U+.47)+.5/(1+\log(T+1))-.46<0\).
In particular \(\Delta>U+.95\) for \(h\ge4\).
For the prime-count term, \(\log P\ge
L/4+\log2+(\log2048)/4>L/4+2.59\).
Write \(x=\log P\) and \(C_\pi=1.2762\).
Then \(L(1+C_\pi/x)/x<4(1-C_\pi^2/x^2)<4\).
Since \(P/H\le2.5\), Lemma~\ref{lem:prime-estimates}
gives \(\pi(P)L/H<10\).
By \eqref{eq:large-two-y}, for \(h\ge4\),
\[
 h<\frac{6(3.24U+1)}
           {U+.95-\log2.5-10-3.24U/8192}<61 .
\]
For \(U\ge15\) the denominator is positive and affine in \(U\),
with positive slope and negative constant term.
The numerator is affine with positive slope and constant term,
so the quotient decreases. Thus \(h\le60\).

We next refine the digit loss in \eqref{eq:pade-norm}.
For \(p>P'\ge2H+d\) and a positive lifted valuation \(e\),
one has \(e\ge1/H\) and \(e-\mu\ge1/(2H)\).
The progression in \eqref{eq:pade-progression} has valuation
at most \(N\mu+\lfloor\log_p(bm+j)\rfloor\), as \(p\nmid d\).
The floor first increases beyond \(L_p=\lfloor\log_p M\rfloor\)
only at \(j-1\ge b>2H\). Later increases are more than \(2H\)
steps apart. The penalty \((j-1)(e-\mu)\) absorbs all these
increases. Thus above \(P'\) the per-prime digit loss in the
averaged bound is at most \(L_p\log p/N\), with \(L_p\le3\).
In particular the digit loss is at most
\begin{equation}\label{eq:large-digit}
 \frac{\pi(P')V}{N}+\frac3N\sum_{p>P'}v_p(F(0))\log p .
\end{equation}

Put \(\delta=H/N\) and take \(P'=(2+2/h)H=rH\).
For nonexceptional primes above \(P'\),
Lemma~\ref{lem:two-heights} puts their entire norm product
in the binomial coefficient
\(\binom{b+H+d-1}{2H+2d-1}\).
With \(c=b/H>6.9\), its logarithm divided by \(H\) is at most
\(r(1+\log((c+r/2)/r))\). Its derivative in \(r\) is
\(\log(c/r+1/2)+r/(2c+r)>0\) for \(r\le3\).
At \(r=3\), divide by \(3\) and write the result as
\(t+1-\log(3h)+\log(1+1.5/c)<t-.595\), using \(h\ge2\). For the primes below \(P'\), the same prime-count
calculation as for \(P\), now with \(V<L+.001\) and
\(\log P'>L/4+2.59\), gives \(\pi(P')V/N<4r\delta\le12\delta\).
The exceptional primes contribute at most \(6\delta L/(Br^3)\).
Therefore the digit loss is less than
\(\delta(9t+J_0(U))\), where \(J_0(U)=6.65+6(3.24U)/16384\).

When \(d\ge50\) the weaker angular bound gives
\(\log\Lambda>U+.98t-1.64\).
Use \(t>(U-\log(2U+2))/2\), \(g_0/d<.025\), and
\((2/\pi)\arcsin(e^{-t/2})g_0<.19\).
For the last bound let \(x=e^{-(U-\log\sigma)/4}\).
The estimate \(\arcsin x\le x/\sqrt{1-x^2}\) bounds the loss
by \(2x(1.266+.256U)/(\pi\sqrt{1-x^2})\).
The logarithmic derivative of \(x\) is at most \(-15/64\),
and that of \(1.266+.256U\) is below \(1/U\).
Also \(1/\sqrt{1-x^2}\) decreases. Thus the maximum is at \(U=15\). Split the prime sum at \(HV>e^8\).
If \(\delta\le .045\), the gain exceeds the two upper bounds
by more than
\[
 (.98-9\delta)t-\log\delta-\log(3.24U)-2.101-J_0(U)\delta.
\]
On \(\delta\le.045\) the expression decreases in \(\delta\)
and increases in \(t\). At \(\delta=.045\) and
\(t=(U-\log\sigma)/2\), its derivative in \(U\) is at least
\[
 \frac{.98-9(.045)}2\left(1-\frac1{16}\right)
           -\frac1{15}-.045\frac{6(3.24)}{16384}>0 .
\]
Its value exceeds \(.10\) at \(U=15\).
For \(d<50\), \(s<2500/N<.001\) and \(H<3000\), so the
digit bound is below
\(3000e^{-U}(9U+J_0(U))\). This decreases for \(U\ge15\)
and is below \(.14\) there. The gain exceeds \(U+t/3-.71\), and \(t>U-\log50\).
It exceeds \(\log(3000\cdot3.24U)+.461+.14\) at \(U=15\) and
increases faster. Therefore
\begin{equation}\label{eq:large-degree}
 h\le60,\qquad \delta>.045,\qquad \Delta_*>U-1.01 .
\end{equation}

\subsection*{A count of terminal square multiples}

We use a sharper count in the range \eqref{eq:large-parameter-range}.
For either \(i=N-1,N\), set \(X=bi\), and count primes \(p>P\)
with \(X-2H\le qp^2\le X\), \(q\) a positive integer.
Divide the range up to \(\sqrt X\) into intervals
\([D,\eta D)\) with \(D=\eta^j P\) and \(\eta=9/8\).
The integer points \((p,q)\) lie between \(X/p^2\) and
\(X/p^2-\epsilon\), where \(\epsilon=2H/D^2\).
Each prime permits at most one multiplier and each multiplier
at most one prime. The mean value theorem gives
\(1\le2X|p_2-p_1|/D^3+\epsilon\) for distinct points.
Since \(\epsilon\le1/8\), their abscissas have spacing at least
\(g_D=\max(2,7D^3/(16X))\).

Partition each interval further into blocks of width
\(\rho=(D^4/X)^{1/3}\), with number at most
\(.132(X/D)^{1/3}\), since \(\rho/D\le X^{-1/6}<.007\).
For three points in a block the integer area determinant
is the determinant for \(X/x^2\) plus that for the vertical
errors. The first is at most \(3/4\) in absolute value by
the second divided difference bound \(3X/D^4\), times
\(\rho^3/4\). The error determinant is at most
\(\epsilon\rho\le(2H/X)^{1/3}<.01\). Thus all points in a
block are collinear. For \(n\ge3\) points, select a middle
one whose horizontal distances to both extremes are at least
\((n-2)g_D/2\). The chord of \(X/x^2\) through the extreme
abscissas is at most \(\epsilon\) above the curve there.
If its two distances are \(L_1,L_2\), the chord excess is at
least \(3XL_1L_2/(\eta D)^4\), by subtracting a quadratic of
second derivative \(6X/(\eta D)^4\) and using convexity.
The upper bound \(\epsilon=2H/D^2\) on this excess gives
\[
 n\le 2+2\eta^2\sqrt{2/3}\,
                 \frac{D\sqrt{H/X}}{g_D}.
\]
The same bound holds for smaller \(n\).

Sum these bounds with the weight \(\log(\eta D)\).
The constant term from each block sums to at most
\[
 7(X/P)^{1/3}(\log P+3.06),
\]
since \(.264/(1-\eta^{-1/3})<7\) and
\(\log\eta/(1-\eta^{-1/3})<3.06\).
For the second term split at \(D_*=(32X/7)^{1/3}\),
where \(g_D\) switches. The powers of \(D\) in the two
geometric sums are \(2/3\) and \(-7/3\).
Their combined unweighted coefficient is bounded using
\[
 .132\eta^2\sqrt{2/3}(32/7)^{2/9}
 \left(\frac1{1-\eta^{-2/3}}+\frac1{1-\eta^{-7/3}}\right)<3.4 .
\]
The first sum has weight at most \(\log(\eta D_*)\).
For the second sum, the weighted summand \(D^{-7/3}\log(\eta D)\)
decreases above \(D_*\), as \(\log(\eta D_*)>3/7\).
Replace the actual starting point by \(D_*\) in every term of
its geometric progression. In that upper-bounding series, the
mean weight is
\(\log(\eta D_*)+(\log\eta)\eta^{-7/3}/(1-\eta^{-7/3})
<(\log X)/3+1\).
Summing over both orders and using \(X\le M\) yields
\begin{equation}\label{eq:large-square-charge}
 \frac{Y}{H}<
 \frac{14(\log P+3.06)}{(P/H)^{1/3}B^{1/3}}
 +\frac{6.8 M^{1/18}}{\sqrt H}(L/3+1).
\end{equation}

\subsection*{The final comparison}

Put \(v=N e^{-t}=sh/\delta>6/5\). Then
\[
 M=\frac{\delta N^3}{hv},\qquad B=hv\delta^3N .
\]
Using \(\log P>U+\log(.09)\), \(L<3.24U\),
and \(\Delta_*>U-1.01\) in \eqref{eq:large-norm}
gives \(Y/(HL)>.05\). Explicitly, a lower bound is
\[
 \frac{U-1.01-\log2.5-
       \dfrac{2.5(3.24U)}{U+\log(.09)}
          \left(1+\dfrac{1.2762}{U+\log(.09)}\right)
       -3.24U/8192}{3.24U}.
\]
After division by \(3.24U\), the \(U\)-dependent negative
terms are multiples of \(1/U\), \(1/(U+\log(.09))\), and
\(1/(U+\log(.09))^2\). Each decreases, so the lower bound
is smallest at \(U=15\), where it exceeds \(.0502\).
By \eqref{eq:large-two-y},
\(Y/(HL)\le 1/(v h(2+1/h)^2\delta^2)\). Hence
\(\delta^2<4/3\), since \(h(2+1/h)^2\ge12.5\).
In particular \(\delta<1.4\) and \(L<3U\).

For \(\ell\le\delta\le u\) and \(v\ge w\), define
\[
\begin{aligned}
 V_0&=3U+\log(u/(2w))+.001,\\
 Z&=\frac1\ell\sum_{j=4}^8(u/(2w))^{1/j}e^{(3/j-1)U}
       +\frac{3U}{\ell\log2}(u/(2w))^{1/9}e^{-2U/3},\\
 C_0&=1.017\{5+\min(2.5,(2w\ell^2)^{-1/3})+Z\},\\
 A(U,\ell,u,w)&=U-1.01-\log2.5-C_0-V_0/8192 .
\end{aligned}
\]
At \(p\le P\), the extra term in the affine norm count is
\(\sum_{j\le L/\log2}\vartheta(\min(P,M^{1/j}))\), where
\(\vartheta(x)=\sum_{p\le x}\log p\).
By Rosser and Schoenfeld,
\(\vartheta(x)<1.017x\) for \(x>0\) and
\(\vartheta(x)>.98x\) for \(x\ge7481\)
\cite[Theorems 9 and 10]{rosser-schoenfeld}.
In the upper estimate, \(j=1,2\) contribute at most
\(1.017(2P/H)\) after division by \(H\), the term \(j=3\)
uses \(M^{1/3}/H\le(2w\ell^2)^{-1/3}\), and the remaining
terms contribute at most \(1.017Z\).
Thus \eqref{eq:large-norm} with this count gives \(Y/H>A\).
On the other hand, \eqref{eq:large-square-charge} gives
\(Y/H<W(U,\ell,w)\), with
\[
 W(U,x,w)=e^{-U/3}\left\{
 \frac{14(U+\log(2.5x)+3.06)}{(5w)^{1/3}x}
 +\frac{6.8(U+\tfrac13\log(x/(2w))+1)}{(2w)^{1/18}x^{4/9}}
 \right\}.
\]
In obtaining this formula first take \(x=\delta\), and use
\(hv\ge2w\), \((P/H)hv\ge5w\), and
\(\log P\le U+\log(2.5\delta)\).
The numerators in braces have the form
\(\varphi=U+c\log x+c'\), with \(c=1,1/3\), respectively,
and the powers of \(x\) in the denominators are
\(\nu=1,4/9\). For every row below,
\(.045\le x\le1.4\), \(6/5\le w\le9\), and both values
of \(\varphi\) exceed \(14\).
The \(x\)-derivative of \(x^{-\nu}\varphi\) has the sign
of \(c-\nu\varphi<0\). After multiplying by \(x\) the sign
is that of \(c+(1-\nu)\varphi>0\).
Thus \(W\) decreases in \(x\) and \(xW\) increases.
In particular \(Y/N<uW(U,u,w)\).
The derivative of \(e^{-U/3}\varphi\) in \(U\)
has the sign of \(1-\varphi/3<0\), so \(W\) decreases.
In \(A\), each term of \(Z\) decreases, including the last
one since \(U>3/2\). Therefore \(A'(U)>1-3/8192>0\).

At \((\ell,u,w)=(.21,1.4,6/5)\) the difference \(A-W\)
exceeds \(.074\) at \(U=15\), so \(\delta<.21\).
On the resulting range \([.045,.21]\) with \(w=6/5\),
\(A/V_0\) increases with \(U\). To check this, write
\(A/V_0=(U-c(U))/(3U+\beta)-1/8192\),
where \(c(U)=1.01+\log2.5+C_0\) decreases and
\(\beta=\log(.21/2.4)+.001\) satisfies \(\beta+3c(U)>0\).
The multiplier bound in \eqref{eq:large-two-y} now gives
\(sh<.675/\delta\), since \((2/25)V_0/A<.675\) at \(U=15\).
For larger \(\delta\), the table below uses the stronger
bound \(sh<uw\) when \(A>W\) excludes \(v\ge w\).

Now \eqref{eq:large-angular-template} gives
\(\log\Lambda>U+t-1.60\). Indeed \((t+g_0)/d<.009\)
from \(d>.045N/60\) and \(t<U\). The angle loss is below
\(.19\) as before, and \(c_\alpha<1.39\) since
\(\alpha<\sqrt{\sigma/N}\).
For each row \([\ell,u]\) put \(V_0=V_0(U,u,6/5)\) and set
\[
 Q=\min\{V_0,\ U+.127+W(U,\ell,6/5)+V_0/8192\}.
\]
Both entries bound the large-prime log norm divided by \(H\).
For the second entry, the equal-height prime contribution
is at most \((2/h)(U+.127)\), since
\(A_0+2<(5/6)N+1/6\), and the other contributions use
\eqref{eq:large-square-charge} and the cube bound.
Both branches of \(Q\) exceed \(15\), so \(HQ>P\)
and \(HQ>e^8\). Below \(HQ\), use \(S(HQ)\).
Above \(HQ\), all primes belong to the large-prime part
and their logarithmic contribution to the prime sum is
at most \(HQ/(HQ-1)<1.001\). This gives \(\log(HQ)+.461\).

The digit estimate \eqref{eq:large-digit} has \(L_p\le2\)
above \(P\), except that \(L_p\le3\) when \(p\le M^{1/3}\).
Since \(P>2H>.09N>7481\), the lower bound for
\(\vartheta(P)\) applies. The extra loss is at most
\(\max(0,\vartheta(M^{1/3})-\vartheta(P))/N\), which is
bounded by \(\max(0,1.017(u/2.4)^{1/3}-1.96\ell)\).
Thus the digit loss is at most the following, with
\(z=U+\log(2.5\ell)\):
\[
\begin{aligned}
 D_0(U,\ell,u)
  &=u\left\{\frac{2.5V_0}{z}\left(1+\frac{1.2762}{z}\right)
                       +2(U+.127)+V_0/4096\right\}\\
  &\qquad+2uW(U,u,6/5)+\max(0,1.017(u/2.4)^{1/3}-1.96\ell).
\end{aligned}
\]
For the prime count in this formula put \(y=\log(rH)>8\).
The derivative in \(r\) of \(r(1+C_\pi/y)/y\) is
\(\bigl(y^2+(C_\pi-1)y-2C_\pi\bigr)/y^3>0\).
Thus the prime-count contribution is bounded using \(r=2.5\). Since \(H=she^t\), the surplus over
\eqref{eq:pade-norm} is greater than
\(U-\log(JQ)-2.061-D_0\) where \(J\) bounds \(sh\).
The following endpoint comparisons complete the proof:
\[
\begin{array}{c|c|r|c|c}
{}[\ell,u]&w&J&
\substack{A-W\\\text{at }U=15}&
\substack{U-\log(JQ)-2.061-D_0\\ \text{at }U=15}\\ \hline
{}[.045,.10]&\text{--}&15&\text{--}&>.36\\
{}[.10,.14]&9&1.26&>.109&>1.63\\
{}[.14,.18]&3.6&.648&>.117&>.83\\
{}[.18,.21]&1.8&.378&>.093&>.30
\end{array}
\]
The first row uses \(sh<.675/\delta\). In the other rows
the fourth column gives \(v<w\). To check the last column throughout the range, first note
\(V_0>40\) in all four rows. The two branches defining \(Q\)
have logarithmic derivatives bounded above by \(3/40\) and
\((1+3/8192)/(15+.127)\), respectively. For the second
bound, use \(W>0\) in the denominator and \(\partial_U W<0\)
in the numerator. Both bounds are below \(.08\).
This bounds the growth of \(\log Q\) even at a change
of the minimizing branch, since each branch minus \(.08U\)
in logarithmic form is decreasing.

The rational factor \(V_0(1+C_\pi/z)/z\) equals
\(3(1-k/z)(1+C_\pi/z)\), where
\(k=\log(2.5\ell)-(\log(u/2.4)+.001)/3<.02\).
Its derivative is \(3(-(C_\pi-k)z+2kC_\pi)/z^3<0\)
for \(z>8\).
The \(W\) term also decreases. The remaining \(U\)-dependent
terms of \(D_0\) grow at rate at most
\(.21(2+3/4096)<.43\). Consequently the last column increases
with slope at least \(1-.08-.43>0\) wherever the derivative exists
and is continuous at a switch of branch.
Thus the positive surplus contradicts the Padé inequality.
The accompanying verifier reproduces the scalar comparisons
at the stated endpoints using rational enclosures for logarithms,
exponentials, and rational powers.
\end{proof}



\begin{thebibliography}{99}

\bibitem{dusart-primes}
P. Dusart,
\emph{Estimates of some functions over primes without R.H.},
\href{https://arxiv.org/abs/1002.0442v1}{arXiv:1002.0442v1}, 2010.

\bibitem{dusart-interval}
P. Dusart,
\emph{Explicit estimates of some functions over primes},
Ramanujan J. \textbf{45} (2018), 227--251.
\href{https://doi.org/10.1007/s11139-016-9839-4}{doi:10.1007/s11139-016-9839-4}.

\bibitem{edo-essen}
E. Edo and A. van den Essen,
\emph{The Strong Factorial Conjecture},
\href{https://arxiv.org/abs/1304.3956v2}{arXiv:1304.3956v2}, 2013.

\bibitem{filaseta-rocks}
M. Filaseta and B. Rocks,
\emph{On the irreducibility of a polynomial associated with the Strong Factorial Conjecture},
Colloq. Math. \textbf{145} (2016), 307--314.
\href{https://doi.org/10.4064/cm6761-3-2016}{doi:10.4064/cm6761-3-2016}.

\bibitem{laishram-shorey}
S. Laishram and T. N. Shorey,
\emph{The greatest prime divisor of a product of consecutive integers},
Acta Arith. \textbf{120} (2005), 299--306.
\href{https://doi.org/10.4064/aa120-3-5}{doi:10.4064/aa120-3-5}.

\bibitem{nguyen-dang}
K.-H. Nguyen-Dang,
\emph{Uniform irreducibility and Galois groups of generalized Laguerre polynomials},
preprint \href{https://hexagonmath.org/pdf/2611.00087v1}{hexagon:2611.00087v1},
2026.

\bibitem{oliveira-gaps}
T. Oliveira e Silva, S. Herzog, and S. Pardi,
\emph{Empirical verification of the even Goldbach conjecture and computation
of prime gaps up to \(4\cdot10^{18}\)},
Math. Comp. \textbf{83} (2014), 2033--2060.
\href{https://doi.org/10.1090/S0025-5718-2013-02787-1}{doi:10.1090/S0025-5718-2013-02787-1}.

\bibitem{rocks}
B. Rocks,
\emph{On the incompatibility of diophantine equations arising from the
Strong Factorial Conjecture},
J. Algebra \textbf{478} (2017), 341--352.
\href{https://doi.org/10.1016/j.jalgebra.2017.01.039}{doi:10.1016/j.jalgebra.2017.01.039}.

\bibitem{rosser-schoenfeld}
J. B. Rosser and L. Schoenfeld,
\emph{Approximate formulas for some functions of prime numbers},
Illinois J. Math. \textbf{6} (1962), 64--94.
\href{https://doi.org/10.1215/ijm/1255631807}{doi:10.1215/ijm/1255631807}.

\bibitem{wagner-stability}
D. G. Wagner,
\emph{Multivariate stable polynomials: theory and applications},
Bull. Amer. Math. Soc. \textbf{48} (2011), 53--84.
\href{https://doi.org/10.1090/S0273-0979-2010-01321-5}{doi:10.1090/S0273-0979-2010-01321-5}.

\end{thebibliography}


\end{document}

